\documentclass[12pt,a4paper]{article}

\usepackage[T1]{fontenc}

\usepackage{amsmath,amsthm,amssymb}
\usepackage{mathtools}
\usepackage{geometry}
\usepackage{microtype}
\usepackage{setspace}
\usepackage{titlesec}
\usepackage{abstract}
\usepackage{hyperref}
\usepackage{enumitem}
\usepackage{booktabs}

\geometry{a4paper,top=2.5cm,bottom=2.5cm,left=3cm,right=3cm}
\onehalfspacing

\hypersetup{colorlinks=true,linkcolor=black,citecolor=black,urlcolor=black}

\newtheorem{theorem}{Theorem}[section]
\newtheorem{lemma}[theorem]{Lemma}
\newtheorem{corollary}[theorem]{Corollary}
\newtheorem{proposition}[theorem]{Proposition}
\newtheorem{definition}[theorem]{Definition}
\theoremstyle{remark}
\newtheorem{principle}{Principle}[section]
\newtheorem{conjecture}{Conjecture}[section]
\newtheorem{remark}{Remark}[section]

\titleformat{\section}{\large\bfseries}{\thesection.}{0.6em}{}
\titleformat{\subsection}{\normalsize\bfseries}{\thesubsection.}{0.5em}{}
\titleformat{\subsubsection}{\normalsize\itshape}{\thesubsubsection.}{0.5em}{}

\title{\textbf{Quadratic Gravity as a Derived Interface Theory:\\
Admissibility Manifolds, Projection Geometry,\\
and the Emergence of Observable Physics}}

\author{Flyxion\\[0.4em]
\small Independent Researcher}

\date{\today}

\begin{document}

\maketitle

\begin{abstract}
We argue that the standard rejection of higher-derivative quantum gravity
on grounds of ghost instability rests on an assumption stronger than
experiment requires: that the internal state space of a theory is
identical to its observable sector. Relaxing this identification, we
develop a framework of admissibility manifolds in which physical
constraints are imposed on projected observable histories rather than
on the coordinates of the ambient state space. Within this framework,
negative-norm sectors need not constitute observational pathologies; they
may instead correspond to hidden admissibility directions that survive
internally while canceling or decoupling under projection to observable
quantities. We formulate the Observable Admissibility Theorem and the
Projection Sufficiency Principle, and show that Krein-space quantum
mechanics---independently proposed by Turok and Bateman for quadratic
gravity---arises as a special case of admissibility-flow dynamics. The
modified Born rule of Turok and Bateman, expressed via projection operators
and a ghost-parity symmetry, is derived from the general principle that
probabilities are computed from projected quantities on $M_{\mathrm{obs}}$
rather than from inner products on $M_{\mathrm{adm}}$. We further argue
that the Born rule itself admits a coarse-graining derivation: following
Needham's geometric reading of complex multiplication as rotation-scaling
transport and Barandes's reconstruction of quantum theory from stochastic
processes admitting a unitary lift, we show that probabilities arise from
projection of field-valued rotation-scaling transport on $M_{\mathrm{adm}}$
onto observable transition structure on $M_{\mathrm{obs}}$; the resulting
transition kernel is unistochastic when the fiber of $\Pi$ admits a unitary
action preserved under the dynamics. We further connect this framework to
Observational--Interventional Separation, the critique of latent
fundamentalism, and the Derived Interface Conjecture: that observable
spacetime theories are effective interfaces generated by projection from a
larger admissibility manifold, with gauge symmetry, Hilbert structure,
metric positivity, and the Born rule itself appearing as interface
properties rather than ontological necessities.
\end{abstract}

\tableofcontents
\newpage

%----------------------------------------------------------------------
\section*{Prefatory Note}
\addcontentsline{toc}{section}{Prefatory Note}
%----------------------------------------------------------------------

The arguments developed here belong to a research program centered on the
conviction that reachability, admissibility, and constraint closure are the
correct ontological primitives---prior to objects, representations,
prediction, and optimization. The present paper applies that orientation to
a problem in the foundations of physics that ramifies far beyond its
motivating example: when is a pathology in the representational machinery
of a physical theory a pathology of the world the theory describes?

The motivating example is provided by Neil Turok and Sam Bateman's proposal
to rehabilitate quadratic gravity using a Krein-space formulation of quantum
mechanics with a modified Born rule based on ghost-parity symmetry. Their
proposal attacks a specific assumption---that the state space of a physical
theory must be a positive-definite Hilbert space---and argues that this
assumption is stronger than observational requirements demand. But the
deeper claim of the present paper is that the Turok--Bateman move is one
instance of a structural distinction that recurs throughout physics and
philosophy of science: the distinction between representation and observation.

That distinction appears in many guises, and the paper traces them
deliberately in parallel. The identification of a gauge orbit with a
physical state, the identification of a wave function with a measurement
outcome, the identification of a Markov blanket with an agent boundary,
the identification of a model coordinate with an ontological object---all
are instances of the same error, which we call latent fundamentalism. The
admissibility framework is a systematic attempt to replace that error with
a principled criterion: physical constraints belong on the projected
observable sector, not on the internal representational machinery.

The paper may therefore be read in two ways. As a contribution to quantum
gravity, it argues that higher-derivative theories with ghost sectors are
not automatically inadmissible, and that quadratic gravity in particular
deserves rehabilitation. As a contribution to the theory of physical
representation, it argues that many historical disputes---about ghosts,
gauge redundancy, wave function realism, and the multiverse---arise from
confusing the geometry of representation with the geometry of observation.
The quadratic gravity discussion provides the most technically demanding
test case; the Observational--Interventional Separation theorem and the
critique of latent fundamentalism are intended to survive it.

The paper proceeds in nine main sections and four appendices. Sections~1--3
establish the framework of admissibility manifolds and prove the two
principal theorems. Sections~4--5 analyze the Ostrogradsky instability and
Krein-space formalism, recovering the Turok--Bateman construction as a
derived consequence. Section~6 addresses the emergence of quantum probability
from projection. Section~7 interprets quadratic gravity as effective
admissibility-flow geometry, with Einstein gravity and quadratic
gravity as successive terms in a curvature expansion. Section~8 develops
the epistemological foundations through Observational--Interventional
Separation and the critique of latent fundamentalism. Section~9 states
and discusses the Derived Interface Conjecture.


%======================================================================
\section{The Hilbert Assumption and Its Hidden Content}
%======================================================================

\subsection{The Orthodox Framework}

Standard quantum mechanics rests on two closely related commitments. The
first is that the state space of a physical system is a complex Hilbert
space $\mathcal{H}$, equipped with a positive-definite inner product
satisfying $\langle\psi\mid\psi\rangle > 0$ for all nonzero $\psi$. The
second, seldom made explicit, is that this space \emph{is} the physical
space: every element of $\mathcal{H}$ is a candidate physical state, and
every self-adjoint operator on $\mathcal{H}$ is a candidate observable.
The Born rule connects them by assigning to each measurement outcome $F$
the probability
\begin{equation}
  P(F \mid I) = |\langle F \mid S \mid I \rangle|^2,
\end{equation}
where $\lvert I\rangle$ and $\lvert F\rangle$ are normalized initial and
final states and $S$ is the $S$-matrix.

When quantum mechanics is extended to incorporate gravity, the second
commitment creates difficulties. The Einstein--Hilbert action
\begin{equation}
  S_{\mathrm{EH}} = \frac{1}{16\pi G}\int d^4x\,\sqrt{-g}\,R
\end{equation}
is not perturbatively renormalizable. Achieving renormalizability requires
adding quadratic curvature terms, as first established by Stelle in 1977:
\begin{equation}
  S = \int d^4x\,\sqrt{-g}
    \!\left(\frac{R}{16\pi G} + \alpha R^2 + \beta C_{\mu\nu\rho\sigma}C^{\mu\nu\rho\sigma}
    \right),
\end{equation}
where $C_{\mu\nu\rho\sigma}$ is the Weyl tensor. The resulting propagator
introduces additional poles corresponding to massive spin-2 degrees of
freedom with opposite-sign kinetic term. In the canonical quantization of
such theories, these states carry negative inner product: they are
\emph{ghosts}. The standard inference follows: negative inner product
implies negative probability via the Born rule; negative probability is
unphysical; ghost-containing theories are unphysical.

Turok has described this inference as appearing ``everywhere in the
literature''---and then argued that it is wrong.

\subsection{The Hidden Assumption Identified}

The standard inference contains two steps. The first is mathematical and
unimpeachable: in a positive-definite Hilbert space, states of negative
inner product do not exist by definition. The second step is physical and
non-trivial: it identifies the inner product on the state space with the
empirical probability function. That identification is an assumption.

We call it the \emph{Hilbert Assumption}:
\begin{equation}
  \text{Observable Physics} = \text{Entire State Space.}
\end{equation}
The Hilbert Assumption is what makes the inner product on the state space
directly interpretable as a probability amplitude. It is satisfied in
every physical situation so far tested. But the question is whether it is
a \emph{necessary} condition for a consistent physical theory or merely
a sufficient one that has been promoted to necessity by historical habit.

Turok's observation is pointed: ``A quantum state is nothing but a label
for a system. Its norm is neither here nor there. You can't observe the
norm of a quantum state.'' What you observe are transition probabilities
between states, and those are constrained to be non-negative and to sum
to unity---a weaker condition than requiring positive norm on every element
of the state space. The present paper takes this observation and embeds it
in a general framework: the Hilbert Assumption is a special case; the
general case allows the state space to be larger than the observable sector.

\subsection{Historical Precedents for a Larger State Space}

The move of allowing a larger internal space than the observable space is
not unprecedented in physics. Two examples are immediately instructive.

Quantum electrodynamics in Lorenz gauge employs an indefinite-metric state
space containing negative-norm longitudinal photon states (the Gupta--Bleuler
construction). Physical states are selected by a subsidiary condition that
projects onto the positive-norm sector. The negative-norm states are
present in the formalism because working in the larger space respects the
manifest Lorentz invariance of the theory; they are absent from any
physical prediction.

The BRST formalism for Yang--Mills theories extends this pattern: one
works in a state space that includes Faddeev--Popov ghost fields obeying
the wrong statistics. Physical states are the cohomology classes of the
BRST operator $Q$: states annihilated by $Q$ modulo states of the form
$Q\lvert\chi\rangle$. Again, the ghost fields exist in $M_{\mathrm{adm}}$
but are absent from $M_{\mathrm{obs}}$.

Turok acknowledges this precedent: physicists are ``actually very used
to working in pseudo-Hilbert spaces. The usual prescription is you think
you're in this big space which has both positive and negative and null
states, and then you perform a projection and you say there's a physical
subspace where everything is positive.'' The novelty of the Krein-space
proposal is not the existence of the larger space but a modification of
what counts as the projection: instead of projecting onto a positive-norm
subspace, one traces over the entire space using a modified probability
functional. From the admissibility perspective, both procedures are
instances of the same principle---constraints on observable quantities
rather than on internal coordinates.


%======================================================================
\section{The Framework of Admissibility Manifolds}
%======================================================================

\subsection{The Three-Level Hierarchy}

Let $M_{\mathrm{adm}}$ denote the \emph{admissibility manifold}: the space
of all field histories consistent with the constraint closure conditions
governing the fundamental dynamics.

\begin{remark}[Mathematical status of $M_{\mathrm{adm}}$]
The admissibility manifold is a space of \emph{histories}, not states. An
element $h \in M_{\mathrm{adm}}$ is a field trajectory over a spacetime
region, not a configuration at a moment of time. This trajectory-centric
choice is deliberate and consequential: it permits the framework to address
Ostrogradsky-type pathologies, which are properties of phase-space
directions rather than of trajectories, and it connects naturally to the
path-integral formulation of quantum field theory in which histories are
the fundamental objects.

The appropriate category for $M_{\mathrm{adm}}$ varies by application.
For perturbative quantum field theory, it is a space of distributional
field configurations, and the relevant topology is that of tempered
distributions. For the Krein-space application of Sections~3--5, it is
a Krein space (an indefinite inner-product space with a fundamental
decomposition $\mathcal{K} = \mathcal{K}_+ \oplus \mathcal{K}_-$). For
the RSVP field-theoretic application of Section~6, it is a space of
field-valued transport operators equipped with a fiber bundle structure
over $M_{\mathrm{obs}}$. The framework does not presuppose any single
category; it specifies instead the properties any admissibility manifold
must satisfy: the existence of the surjective projection $\Pi$ and the
well-definedness of the admissibility conditions of Definition~2.1 on its
image. Readers seeking a canonical formal setting may take $M_{\mathrm{adm}}$
to be a Hilbert manifold modeled on a separable Hilbert space, with $\Pi$
a smooth surjective submersion; all results in Sections~2--3 hold in
this setting.
\end{remark}

The geometry of $M_{\mathrm{adm}}$ is not postulated; it is whatever
structure is imposed by the constraint closure conditions of the theory.

Above $M_{\mathrm{adm}}$ sits the \emph{projection manifold}
$M_{\mathrm{proj}}$: equivalence classes of admissible histories under
the symmetries generating non-observable degrees of freedom. Gauge
transformations, diffeomorphisms, BRST-exact states, and ghost-parity
equivalences all generate such classes. Above that sits the
\emph{observable manifold} $M_{\mathrm{obs}}$: empirically
distinguishable outcomes. These spaces are connected by surjective
projections:
\begin{equation}
  M_{\mathrm{adm}}
  \;\xrightarrow{\;\Pi_1\;}\;
  M_{\mathrm{proj}}
  \;\xrightarrow{\;\Pi_2\;}\;
  M_{\mathrm{obs}},
\end{equation}
with composite projection $\Pi = \Pi_2 \circ \Pi_1 : M_{\mathrm{adm}} \to M_{\mathrm{obs}}$.

The Hilbert Assumption corresponds to $\Pi$ being an isomorphism. The BRST
approach treats $\Pi$ as a proper surjection with nontrivial kernel but
requires $M_{\mathrm{adm}}$ to be a Hilbert space. The Krein-space approach
relaxes the latter: $M_{\mathrm{adm}}$ need not be a Hilbert space, provided
$\Pi$ maps it to a space of admissible probability distributions.

\begin{definition}[Observable Admissibility]
\looseness=-1 A history $h \in M_{\mathrm{adm}}$ is \emph{observable-admissible}
if its projection $\Pi(h) \in M_{\mathrm{obs}}$ satisfies:
\begin{enumerate}[label=(\roman*)]
  \item probabilities are non-negative: $P(O) \geq 0$ for all\newline\hspace*{2em}observables $O$;
  \item probabilities for any exhaustive partition sum to unity:
        $\sum_i P(O_i) = 1$;
  \item the causal order on $M_{\mathrm{obs}}$ is preserved under dynamical
        evolution.
\end{enumerate}
We write $\mathcal{A}_O(\Pi(h)) = 1$ when these hold. The admissibility
set in $M_{\mathrm{adm}}$ is
\begin{equation}
  \mathcal{A}_H = \Pi^{-1}\!\bigl(\{\Pi(h) : \mathcal{A}_O(\Pi(h))=1\}\bigr).
\end{equation}
\end{definition}

This definition is deliberately minimal. It does not require the inner
product on $M_{\mathrm{adm}}$ to be positive definite. It does not require
a canonical symplectic structure. It requires only that the image of $h$
under $\Pi$ support a well-defined causal probability measure over
observable outcomes.

\begin{principle}[Projection Sufficiency]
Physical constraints are conditions on projected observable histories.
They need not be conditions on the internal coordinates of $M_{\mathrm{adm}}$.
\end{principle}

Hilbert-space quantum mechanics is recovered as the special case where
$M_{\mathrm{adm}} = \mathcal{H}$ with positive inner product and $\Pi$ is
the identity. The principle allows $M_{\mathrm{adm}}$ to carry richer
structure---including negative-norm directions---provided those directions
do not compromise admissibility under $\Pi$.

\subsection{The Coordinate--Observable Distinction}

A central distinction runs through all that follows. A \emph{coordinate
pathology} is a feature of the internal representation in $M_{\mathrm{adm}}$
with no direct bearing on observable quantities: negative inner product,
gauge redundancy, coordinate singularities, Ostrogradsky's unbounded
phase-space directions. An \emph{observable pathology} is a violation of
one of the three conditions in Definition~2.1: negative probabilities for
measurement outcomes, failure of probability conservation, acausal signaling.

The failure to maintain this distinction accounts for much of the resistance
to higher-derivative theories: coordinate pathologies have been treated as
automatically implying observable pathologies. The following theorem shows
that this is not a general principle.

\begin{theorem}[Projection Pathology Separation]
A coordinate pathology in $M_{\mathrm{adm}}$ does not in general imply an
observable pathology in $M_{\mathrm{obs}}$.
\end{theorem}

\begin{remark}[Logical status of this theorem]
The theorem is a structural consequence of the definitions rather than a
deep result: admissibility is placed on the image of $\Pi$, so whether a
property of the domain induces a property of the image depends on $\Pi$
alone. The theorem's importance is therefore diagnostic rather than
constructive. It identifies the \emph{correct question}---does $\Pi$ map
this pathology to $M_{\mathrm{obs}}$ or to the fiber?---and establishes
that this question cannot be answered by inspecting the pathology alone.
Concretely: showing that a theory has a ghost in $M_{\mathrm{adm}}$ does
not constitute a proof that the theory has an observable instability.
An additional argument about the structure of $\Pi$ is always required.
The historical errors the theorem diagnoses arose from omitting that step.
\end{remark}

\begin{proof}
A coordinate pathology is a property $P$ of some $h \in M_{\mathrm{adm}}$.
Observable admissibility is determined entirely by $\mathcal{A}_O(\Pi(h))$.
Since $\Pi$ is many-to-one, elements of $M_{\mathrm{adm}}$ carrying $P$
may map to elements of $M_{\mathrm{obs}}$ satisfying all three admissibility
conditions. Whether they do so depends on the structure of $\Pi$, not on
$P$ alone.
\end{proof}

The inference from ``this theory has ghost states'' to ``this theory
predicts negative probabilities'' requires an additional premise: that
the projection $\Pi$ does not cancel or absorb the ghost contribution
before reaching $M_{\mathrm{obs}}$. The history of gauge theories
of quadratic gravity.

\paragraph{Interface violation is not repair warrant.}
A terminological caution belongs here. That an interface detects a
projection failure --- a transition admissible in $M_{\mathrm{adm}}$
that fails to remain admissible in $M_{\mathrm{proj}}$ or
$M_{\mathrm{obs}}$ --- identifies a fault. It does not by itself warrant
altering the upstream representation to correct it. Interface violation
and repair warrant are distinct questions: the first is diagnostic,
asking whether admissibility is preserved across projection; the second
is decision-theoretic, asking whether some available intervention
improves the state by more than its cost, given the evidence presently
available for distinguishing improvement from noise.

This distinction is itself interface-relative. A repair operator
situated at $M_{\mathrm{obs}}$ never evaluates some omniscient
$\Delta_{\mathrm{dep}}^{\mathrm{adm}}$; it evaluates a projected
quantity $\Delta_{\mathrm{dep}}^{\mathrm{obs}}(a \mid h,
\mathcal{R}_\Gamma)$ built from distinctions available through
$M_{\mathrm{obs}}$ alone. The interface distance $d_{\mathrm{int}}$
developed above then measures how much intervention-relevant structure
is lost under projection --- how far a judgment made at
$M_{\mathrm{obs}}$ can diverge from the judgment that would be made with
access to $M_{\mathrm{adm}}$. This strengthens rather than complicates
the Observational--Interventional Separation theorem: it is consistent
for $J_{\mathrm{obs}}(a) < 0$ to hold at the observable interface while
the intervention's inaccessible admissibility-level effect is positive,
or vice versa. That is not a defect in the repair objective; it is the
same epistemic limitation the separation theorem already describes,
applied to intervention judgments rather than to static observational
equivalence classes.

The full repair-warrant apparatus --- the reference class
$\mathcal{R}_\Gamma(h)$, the marginal effect $\Delta_{\mathrm{dep}}$,
the objective $J$, and the resulting three-way refusal taxonomy --- is
developed in ``Diagnostic Coordinates and the Time of Repair.'' This
essay's projection results supply the diagnostic and interface-relativity
half of that picture; they do not by themselves decide when correction
is warranted.

%======================================================================

\section{Ghosts as Hidden Admissibility Directions}
%======================================================================

\subsection{The Origin of Ghosts in Higher-Derivative Theories}

In quantum field theory, a ghost is a propagating degree of freedom whose
kinetic term has the wrong sign. For a higher-derivative theory, this arises
from partial-fraction decomposition of the propagator. For quadratic gravity,
the spin-2 part of the propagator takes the schematic form
\begin{equation}
  \Delta^{(2)}(k) \sim \frac{1}{k^2} - \frac{1}{k^2 + m_2^2},
\end{equation}
where $m_2 = M_{\mathrm{Pl}}/\sqrt{2\beta}$. The negative sign on the second
term indicates that the corresponding state has negative inner product with
itself. In canonical quantization, the creation and annihilation operators
for this mode satisfy commutation relations with opposite sign, yielding a
state space with indefinite inner product: the Krein space.

The standard response---declare such states unphysical and seek a projection
onto a positive Hilbert subspace---is unavailable here because no such
projection preserving all the symmetries of quadratic gravity has been found.
Turok and Bateman's response is different: rather than projecting out the
ghost states before computing probabilities, they include them in a modified
probability functional that is engineered to produce non-negative results
regardless of the sign of individual state norms.

\begin{definition}[Hidden Admissibility Direction]
An element $g \in M_{\mathrm{adm}}$ is a \emph{hidden admissibility
direction} if $[g, g]_{\mathcal{K}} < 0$ (a coordinate pathology) and
yet $\mathcal{A}_O(\Pi(g)) = 1$ (no observable pathology).
\end{definition}

\subsection{The Krein Space and Ghost Parity}

A Krein space $\mathcal{K}$ is an indefinite inner product space with
decomposition
\begin{equation}
  \mathcal{K} = \mathcal{K}_+ \oplus \mathcal{K}_-,
\end{equation}
where $\mathcal{K}_+$ carries a positive-definite inner product and
$\mathcal{K}_-$ carries a negative-definite one. The Krein inner product is
\begin{equation}
  [u, v]_{\mathcal{K}} = \langle u_+, v_+\rangle_+ - \langle u_-, v_-\rangle_-.
\end{equation}
This inner product is non-degenerate but indefinite. Turok's Minkowski
analogy is exact: positive-norm states correspond to timelike vectors,
negative-norm states to spacelike vectors, null-norm states to lightlike
vectors. Minkowski spacetime is not deficient as a geometry because
some vectors have negative norm; a Krein space is not deficient as an
algebra because some states have negative inner product.

The \emph{ghost-parity operator} is
\begin{equation}
  G = \Pi_+ - \Pi_-,
\end{equation}
with $\Pi_\pm$ orthogonal projectors onto $\mathcal{K}_\pm$, satisfying
$G^2 = I$ and $G^\dagger = G$. When $G$ commutes with the Hamiltonian
and $S$-matrix, positive-norm and negative-norm sectors evolve independently.
This selection rule is the key to the modified Born rule.

\subsection{The Modified Born Rule and its Derivation from Projection}

In standard quantum mechanics, the Born rule
$P(F\mid I) = |\langle F\mid S\mid I\rangle|^2$ can be rewritten without
explicit normalization as
\begin{equation}
  P(F\mid I) = \mathrm{tr}(\Pi_F\,S\,\Pi_I\,S^\dagger),
\end{equation}
where $\Pi_I = \lvert I\rangle\langle I\rvert / \langle I\mid I\rangle$
and $\Pi_F = \lvert F\rangle\langle F\rvert / \langle F\mid F\rangle$ are
projectors and the trace is over the Hilbert space. This form is
equivalent to the standard Born rule when all state norms are positive but
is well-defined as a functional of projectors even when the inner product
on the full state space is indefinite.

Turok and Bateman adopt exactly this form: defining $A = \Pi_F \cdot S \cdot \Pi_I$,
the probability is $P(F\mid I) = \mathrm{tr}(A^\dagger A)$, where the trace
runs over all states in $\mathcal{K}$, including ghost states. Turok
describes this as ``more economical'' than the traditional BRST projection:
there is no need to construct a physical subspace; one simply sums over
everything and the ghost-parity symmetry ensures the result is positive.

\begin{theorem}[Observable Admissibility]
Let $\mathcal{K} = \mathcal{K}_+ \oplus \mathcal{K}_-$ be a Krein space with
ghost-parity operator $G$. Let the probability functional be
\begin{equation}
  P(F\mid I) = \mathrm{tr}_{\mathcal{K}}(A^\dagger A),
  \quad A = \Pi_F \cdot S \cdot \Pi_I.
\end{equation}
If $G$ is a symmetry of $S$ (i.e., $[G, S] = 0$), then $P(F\mid I) \geq 0$
for all $I, F$, and $\sum_F P(F\mid I) = 1$.
\end{theorem}

\begin{proof}
We proceed carefully to avoid importing Hilbert-space assumptions into the
Krein-space setting.

\textit{Positivity.} The key subtlety is that the adjoint $A^\dagger$ and
the norm $\|A|\psi\rangle\|^2$ must be interpreted with respect to the
\emph{positive} Hilbert space structure on $\mathcal{K}_+$, not the
indefinite Krein inner product. Equip $\mathcal{K}$ with the positive
definite inner product $\langle u, v\rangle_+ = \langle u_+, v_+\rangle +
\langle u_-, v_-\rangle$, where $u = u_+ + u_-$ with $u_\pm \in \mathcal{K}_\pm$.
With respect to $\langle\cdot,\cdot\rangle_+$, the operator $A^\dagger_+ A$
is positive semidefinite, and its trace satisfies
$\mathrm{tr}_+(A^\dagger_+ A) = \sum_n \langle e_n | A^\dagger_+ A | e_n\rangle_+ \geq 0$
for any orthonormal basis $\{|e_n\rangle\}$ of $(\mathcal{K}, \langle\cdot,\cdot\rangle_+)$.
The probability functional is then $P(F\mid I) = \mathrm{tr}_+(A^\dagger_+ A) \geq 0$.

\textit{Normalisation.} We must show $\sum_F \mathrm{tr}_+(A_F^\dagger A_F) = 1$,
where $A_F = \Pi_F S \Pi_I$ and the sum runs over a complete set of
orthogonal final projectors $\{\Pi_F\}$ summing to the identity $I_{\mathcal{K}}$.
Using ghost-parity symmetry $[G, S] = 0$, the $S$-matrix block-diagonalises
with respect to the $\mathcal{K}_\pm$ decomposition:
$S = S_+ \oplus S_-$, where $S_\pm : \mathcal{K}_\pm \to \mathcal{K}_\pm$ are
unitary on their respective sectors. The completeness sum then separates:
\begin{align*}
  \sum_F \mathrm{tr}_+(A_F^\dagger A_F)
  &= \mathrm{tr}_+\!\Bigl(\Pi_I S^\dagger \sum_F \Pi_F S \Pi_I\Bigr)
   = \mathrm{tr}_+\!\bigl(\Pi_I S^\dagger S \Pi_I\bigr).
\end{align*}
Since $S_+$ is unitary on $\mathcal{K}_+$ and $S_-$ is unitary on $\mathcal{K}_-$,
we have $S^\dagger S = I_{\mathcal{K}}$, and the expression reduces to
$\mathrm{tr}_+(\Pi_I) = 1$ (the projector $\Pi_I$ projects onto a single
normalized state). This completes the proof.

The ghost-parity assumption $[G,S] = 0$ is precisely the condition that
keeps $\mathcal{K}_+$ and $\mathcal{K}_-$ from mixing under time evolution.
If this condition fails---if the dynamics couple positive- and negative-norm
sectors---then $S$ does not block-diagonalise, $S^\dagger S \neq I_\mathcal{K}$
in the relevant sense, and normalisation can fail. The ghost-parity
condition is therefore not merely a technical convenience but the essential
physical constraint whose verification in the full spin-2 sector of
quadratic gravity remains open.
\end{proof}

The ghost states in $\mathcal{K}_-$ are hidden admissibility directions:
they participate in the computation of $P(F\mid I)$ through the trace but
do not appear in the final observable probability as negative contributions.
The admissibility interpretation is direct: the trace over $\mathcal{K}$
is a computation on $M_{\mathrm{adm}}$ that produces a quantity on
$M_{\mathrm{obs}}$.

Turok notes that in this construction, the amplitude $A$ itself is not
a physical quantity---``it's not in quantum mechanics, you've got to square
it.'' This observation is significant: it separates the ontological status
of amplitudes (elements of $M_{\mathrm{adm}}$) from that of probabilities
(elements of $M_{\mathrm{obs}}$). The admissibility framework makes this
separation explicit.


%======================================================================
\section{The Ostrogradsky Theorem Revisited}
%======================================================================

\subsection{The Classical Statement}

Ostrogradsky's 1850 theorem establishes that a Lagrangian with
irreducible dependence on derivatives of order higher than the first
possesses a Hamiltonian unbounded below. For a system with Lagrangian
$L(q,\dot{q},\ddot{q})$ in which $\ddot{q}$ appears non-degenerately,
the Ostrogradsky construction introduces canonical coordinates
\begin{align}
  Q_1 &= q, \quad
  P_1 = \frac{\partial L}{\partial\dot{q}} - \frac{d}{dt}\frac{\partial L}{\partial\ddot{q}},\\
  Q_2 &= \dot{q}, \quad
  P_2 = \frac{\partial L}{\partial\ddot{q}},
\end{align}
yielding the Hamiltonian
\begin{equation}
  H = P_1 Q_2 + P_2\,\ddot{q}(Q_2, P_2) - L,
\end{equation}
which is linear in $P_1$ and hence unbounded below. The standard inference
is that coupling such a system to any conventional positive-energy system
produces runaway energy exchange: the higher-derivative system lowers its
energy without bound while the coupled system raises its energy without
bound. The theorem is mathematically correct as a statement about the
extended phase space $M_{\mathrm{adm}}^{\mathrm{ext}}$.

\subsection{Gravity's Native Negative Energy and the Trajectory Shift}

Turok's central observation is that gravity already involves negative
energy in a way that no one regards as problematic: the gravitational
potential energy of a bound system is negative, and the total energy of
an isolated gravitating system can be less than that of its widely
separated constituents. More strikingly, observations show that the
universe is expanding exponentially---a process that, out of context,
resembles an instability. ``That sounds awfully like an instability,''
Turok remarks, ``but it is absolutely stable'' when analyzed as a
cosmological solution.

When Turok and Bateman analyze the four-derivative gravitational theory,
they find that the Ostrogradsky runaway direction in the extended phase
space corresponds precisely to normal exponential cosmological expansion.
The instability is not an observable pathology; it is a coordinate
description of ordinary cosmological behavior.

This is the trajectory shift. State-based reasoning asks: what is the sign
of the energy of this configuration? Trajectory-based reasoning asks: what
are the actual histories generated by the dynamics, and are they admissible
in $M_{\mathrm{obs}}$? The Ostrogradsky theorem establishes a property of
the state space; it does not establish a property of the generated histories.

\begin{principle}[Observable Stability]
Observable stability is determined by the admissibility of projected
trajectories, not by the sign of the Hamiltonian on $M_{\mathrm{adm}}$.
\end{principle}

The Pais--Uhlenbeck oscillator provides a canonical illustration. Its
Hamiltonian is linear in one canonical momentum and hence unbounded below.
Yet its equations of motion reduce to those of two coupled simple harmonic
oscillators with frequencies $\omega_1$ and $\omega_2$: all solutions
are bounded and periodic. The unbounded directions in the extended phase
space are genuine, but no admissible trajectory reaches them. This is
exactly the Projection Pathology Separation Theorem in concrete form: the
coordinate pathology (unbounded Hamiltonian) does not propagate to an
observable pathology (actual runaway behavior) because the dynamics
confines admissible trajectories to the bounded sector of $M_{\mathrm{adm}}^{\mathrm{ext}}$.

\subsection{From State Properties to History Properties}

A recurrent pattern in the admissibility program is the replacement of
state-property reasoning with history-property reasoning. This replacement
arises because observable physics is inherently historical: what we measure
are outcomes of processes extended in time, not instantaneous snapshots
of a state space.

For the Ostrogradsky instability, the frozen distinction is between the sign
of the Hamiltonian (a state property) and the boundedness of the actual
trajectories (a history property). Treating the former as a certificate of
the latter is precisely the conflation of coordinate and observable pathologies
that the Projection Pathology Separation Theorem forbids.

Turok's resolution of the Ostrogradsky instability in the gravitational
context is a concrete instance of this replacement: he does not modify the
Hamiltonian, nor does he add a constraint to restrict the phase space. He
instead analyzes the actual trajectories of the theory and finds that they
correspond to known, stable cosmological solutions. The extended phase space
$M_{\mathrm{adm}}^{\mathrm{ext}}$ has unbounded directions; the admissible
trajectories do not reach them.


%======================================================================
\section{Asymptotic Freedom, the Hierarchy Problem, and UV Completeness}
%======================================================================

\subsection{The Renormalizability and Asymptotic Freedom of Quadratic Gravity}

Two properties of the quadratic gravity theory make it particularly attractive
as a candidate quantum theory of gravity, and both are consequences of
the four-derivative structure of the action.

The first is \emph{perturbative renormalizability}. The four-derivative
kinetic terms suppress high-frequency fluctuations more strongly than the
two-derivative Einstein--Hilbert term, and all ultraviolet divergences in
perturbation theory can be absorbed into redefinitions of the coupling
constants $\alpha$ and $\beta$. This was established by Stelle in 1977 and
remains uncontroversial.

The second is \emph{asymptotic freedom}. As shown by Avramidi in the
1980s, the renormalization group running of the couplings causes them to
flow to zero at high energies, analogously to the strong coupling constant
of QCD. At short distances---above the Planck scale---the theory becomes
arbitrarily weakly coupled, reducing to ``just waves which don't interact,''
as Turok describes it. The theory is therefore UV-complete in the same
sense that QCD is UV-complete: it has a well-defined continuum limit.

The analogy with QCD is structurally deep. QCD is asymptotically free in
the UV and strongly coupled in the IR, with confinement at low energies
ensuring that the only observable excitations are color-neutral hadrons
rather than individual gluons. Quadratic gravity may have a similar infrared
structure: weakly coupled at the Planck scale, strongly coupled at some
lower scale, with a non-trivial spectrum of bound states. Turok notes that
a student is currently investigating this question by studying the theory
on the lattice---the same technique used to study the non-perturbative
IR behavior of QCD.

\subsection{The Spectrum of Excitations}

The physical spectrum of quadratic gravity in the Weyl sector can be
analyzed by decomposing metric perturbations around flat spacetime.
The theory contains four classes of excitations:
\begin{enumerate}[label=(\roman*)]
  \item the standard massless graviton (spin-2, positive norm);
  \item a massive ghost graviton with mass $m_2 = M_{\mathrm{Pl}}/\sqrt{2\beta}$
        (spin-2, negative norm in the Krein space);
  \item a massive vector mode;
  \item a scalar conformal mode governed by the Ricci scalar term.
\end{enumerate}
The scalar mode is the subject of Turok and Bateman's current results: in
the limit $\beta \to 0$ (Weyl coupling goes to zero), the spin-2 ghost
and vector modes decouple and the remaining theory is the scalar conformal
sector. They show that this sector admits the ghost-parity symmetry and
that the modified Born rule produces positive probabilities. The scalar
sector is, in Turok's words, ``a kind of toy model for quantum gravity''---
not the full theory, but a limit of it in which the construction can be
verified rigorously.

The open question is whether the ghost-parity symmetry extends to the full
theory with the spin-2 sector included. Turok acknowledges: ``The trick we
used to make sense of it may or may not apply to the full thing. We will
have to search in the full theory: is there a similar discrete symmetry?
If there is, then this will be a complete theory of quantum gravity.''

\subsection{The Hierarchy Problem and Dimensional Transmutation}

One of the deepest puzzles in theoretical physics is the hierarchy of mass
scales: the Planck mass at $10^{19}\ \mathrm{GeV}$, the weak scale at
$\sim 100\ \mathrm{GeV}$, the QCD scale at $\sim 1\ \mathrm{GeV}$, and
the cosmological constant at $\sim 10^{-3}\ \mathrm{eV}$. Why are these
scales so widely separated? In the standard model, the Higgs mass must be
tuned to prevent radiative corrections from pulling it up to the Planck
scale.

Turok notes that quadratic gravity may resolve this through dimensional
transmutation. An asymptotically free theory defined at the Planck scale
with a small coupling constant will develop a strongly-coupled scale
exponentially below the Planck scale, by the same mechanism that makes
the QCD scale exponentially below the Planck scale without any fine-tuning.
If the Higgs is a composite of the scalar conformal mode of quadratic
gravity, its mass is naturally set by this dynamically generated scale
rather than by a direct coupling to the Planck scale.

From the admissibility perspective, dimensional transmutation is a statement
about the structure of $M_{\mathrm{adm}}$ in the infrared: the admissible
trajectory space of the theory has a non-trivial structure at the dynamically
generated scale that is invisible in the UV description.


%======================================================================
\section{Quantum Probability as Projection}
%======================================================================

\subsection{The Unistochastic Structure}

The Turok--Bateman reformulation of the Born rule connects to a picture in
which quantum probability is not a foundational axiom but an emergent
structure arising from coarse-graining over unobserved degrees of freedom.

Consider a field-theoretic dynamics on $M_{\mathrm{adm}}$ in which the
fundamental degrees of freedom are field configurations $(\Phi,\mathbf{v},S)$
subject to coupled constraint equations. Observable transition probabilities
arise from integrating over all trajectories connecting initial to final
configurations, weighted by the squared modulus of a field amplitude:
\begin{equation}
  T_{ij} = \int_{\gamma : C_i \to C_j} |\Phi|^2\,d^4x.
\end{equation}
The matrix $T$ with entries $T_{ij}$ is doubly stochastic: non-negative
entries (squares of moduli) summing to unity in each row and column. This
is a \emph{unistochastic} matrix.

The Born rule of standard quantum mechanics is the special case in which
the integral over trajectories reduces to a sum over paths between
eigenstates and $\Phi$ reduces to the quantum-mechanical amplitude $\psi$.
The full field-theoretic expression is more general: it includes contributions
from all degrees of freedom in $M_{\mathrm{adm}}$, whether or not they are
directly observable.

\subsection{The Born Rule as Derived Quantity}

Under this interpretation, the Born rule is a derived quantity. It is the
leading-order approximation obtained when the coarse-graining projection
$\Pi$ is nearly trivial---when unobserved degrees of freedom decouple from
observed ones so cleanly that the full integral over $M_{\mathrm{adm}}$
reduces to the modulus-squared of a Hilbert-space amplitude.

The Turok--Bateman modified Born rule $P(F\mid I) = \mathrm{tr}(A^\dagger A)$
is not a revision of quantum mechanics; it is a more faithful expression of
the same underlying probability computation for the case in which the
decoupling is less complete. The ghost sector cannot be simply discarded; it
must be integrated over. The trace in $\mathrm{tr}(A^\dagger A)$ is exactly
this integration: summing over the unobserved degrees of freedom---including
the ghost sector---while producing an observable quantity. Ghost states in
$\mathcal{K}_-$ are trajectories in $M_{\mathrm{adm}}$ that pass through
negative-norm configurations; their contribution to $T_{ij}$ is included,
but their individual norms play no direct role in the observable probability.

\begin{remark}
Turok's observation that the amplitude $A$ ``is not a physical thing''
connects to the admissibility distinction between $M_{\mathrm{adm}}$ and
$M_{\mathrm{obs}}$. Amplitudes are elements of $M_{\mathrm{adm}}$;
probabilities are elements of $M_{\mathrm{obs}}$. The squaring operation
in the Born rule is the projection $\Pi$.
\end{remark}

\subsection{Unsquared Numbers and Unistochastic Coarse-Graining}

The argument so far has been algebraic: we have shown that probabilities
arise from traces of positive semidefinite operators, and that ghost states
contribute to the trace without violating positivity. But there is a more
geometric way to see why probabilities are squares of something deeper,
and why the ``something deeper'' includes directions---phases, rotations,
negative norms---that are not themselves observable. This geometric picture
connects the Krein-space construction to a broader principle about the
structure of measurement.

\subsubsection*{The Unsquared Viewpoint}

In Needham's geometric reading of complex analysis \cite{needham1997},
a complex number $z = re^{i\theta}$ is not merely a number but a compound
operator on the plane: it scales by $r$ and rotates by $\theta$. The imaginary unit $i$
is therefore an ``unsquared'' direction, not because it is mysterious, but
because it encodes a quarter-turn that the algebraic relation $i^2 = -1$
hides behind a sign. Multiplying two complex numbers compounds their
scalings and adds their angles; taking the modulus squared discards the
angle and retains only the scaling:
\begin{equation}
  |z_1 z_2|^2 = |z_1|^2 |z_2|^2, \qquad
  \arg(z_1 z_2) = \arg z_1 + \arg z_2\;[\text{lost under } |\cdot|^2].
\end{equation}
The Born rule $P = |\langle F|S|I\rangle|^2$ is exactly this operation:
it discards the phase accumulated during evolution and retains only the
amplitude ratio. The probability is the squared modulus of an unsquared
rotation-scaling process. Needham's observation that multiplication composes
rotations and scalings makes the Born rule geometrically intuitive: squaring
removes information about orientation while retaining information about
magnitude \cite{needham1997}. What is squared away is the phase---the record
of how much the state turned in Hilbert space during the transition.

The question the admissibility framework asks is: how much more can be
squared away? If the internal space $M_{\mathrm{adm}}$ is larger than
a Hilbert space---if it contains not just phases but also directions of
indefinite norm, gauge redundancies, and ghost sectors---then the squaring
operation (projection $\Pi$) discards correspondingly more. Positive
probabilities can survive because they depend only on what $\Pi$ retains,
not on what it discards.

\subsubsection*{The RSVP Field Operator}

In the RSVP framework, the analogue of the complex operator $z = re^{i\theta}$
is a field-valued operator
\begin{equation}
  Z(x) = \Phi(x)\,e^{\,\widehat{L}(x)},
  \qquad
  \widehat{L}(x) = \mathbf{v}(x)\cdot\nabla + \theta(x)\,\widehat{T},
\end{equation}
where $\Phi(x)$ is the scalar field supplying the scaling component,
$\mathbf{v}(x)$ is the vector field supplying transport along field-space
directions, and $\theta(x)\,\widehat{T}$ is a torsional or phase-like
component encoding rotational structure in the field configuration space.

The operator $Z(x)$ acts on field configurations by simultaneously scaling
their amplitude and transporting them along the vector flow with a torsional
twist. It is the field-theoretic unsquared object: $Z(x)$ itself is not
observable, but $|Z(x)|^2 = |\Phi(x)|^2$ is the local intensity that enters
observable transition weights. The vector and torsional components are
discarded when the modulus is taken---they are the field-space analogue of
the phase in the complex number $re^{i\theta}$.

The trajectory integral over $M_{\mathrm{adm}}$ connecting macrostate $C_i$
to macrostate $C_j$ is
\begin{equation}
  T_{ij}^{\mathrm{RSVP}}
  = \int_{\gamma : C_i \to C_j}
    \bigl|Z[\gamma]\bigr|^2\,\mathcal{D}\gamma,
\end{equation}
where $Z[\gamma]$ is the path-ordered product of $Z(x)$ along the field
trajectory $\gamma$, and $\mathcal{D}\gamma$ is the appropriate measure
on $M_{\mathrm{adm}}$. The transport and torsion components accumulate
and interfere along $\gamma$; their accumulated effect on $|Z[\gamma]|^2$
is the only thing that survives to $M_{\mathrm{obs}}$.

\subsubsection*{Unistochasticity from Phase-Liftability}

\looseness=-1 When the coarse-grained transition kernel $T_{ij}^{\mathrm{RSVP}}$
satisfies a \emph{phase-liftability} condition---when the interference
structure of the $Z[\gamma]$ contributions can be encoded by a unitary
matrix---there exists a unitary $U$ such that
\begin{equation}
  T_{ij}^{\mathrm{RSVP}} = |U_{ij}|^2.
  \label{eq:unistochastic}
\end{equation}
A doubly stochastic matrix expressible in this form is called
\emph{unistochastic}. Birkhoff's theorem establishes that every doubly
stochastic matrix is a convex combination of permutation matrices; the
stronger condition of unistochasticity says it is also the element-wise
square of a unitary matrix, which is a strictly smaller set.

\begin{conjecture}[Projection-Induced Unistochasticity]
Let $\Pi : M_{\mathrm{adm}} \to M_{\mathrm{obs}}$ be a coarse-graining
map whose fibers admit a unitary action preserved under the dynamics
generated by $Z(x) = \Phi(x)e^{\widehat{L}(x)}$. Then the induced
transition kernel on $M_{\mathrm{obs}}$ is unistochastic.
\end{conjecture}

\begin{remark}
This is stated as a conjecture rather than a theorem because the
conditions on the fibers---$\mathrm{U}(n)$-homogeneity and
Haar-measure regularity---are not guaranteed by the general admissibility
framework but must be verified case by case. The proof sketch below
establishes the result under these conditions; extending it to the
full RSVP dynamics without homogeneity assumptions is an open problem.
The conjecture is nevertheless the conceptual hinge of Section~6: it
is the precise statement that the Born rule is a coarse-graining theorem
rather than a foundational axiom.
\end{remark}

\begin{proof}[Proof sketch (under homogeneity assumption)]
Let $\{C_i\}$ be a partition of $M_{\mathrm{obs}}$ into macrostates, and
let $\mathcal{F}_i = \Pi^{-1}(C_i)$ be the corresponding fiber in
$M_{\mathrm{adm}}$. By hypothesis, each fiber $\mathcal{F}_i$ is a
$\mathrm{U}(n)$-homogeneous space. The dynamics generated by $Z(x)$
transports an initial fiber $\mathcal{F}_i$ to a superposition of
terminal fibers $\{\mathcal{F}_j\}$ via a unitary operator $U$ on
$\mathbb{C}^n$ satisfying $Z[\gamma] \mapsto U$ under the fiber action.
The transition weight is
\begin{equation*}
  T_{ij}^{\mathrm{RSVP}}
  = \int_{\gamma : C_i \to C_j} |Z[\gamma]|^2\,\mathcal{D}\gamma
  = \int_{\mathcal{F}_i}\!\int_{\mathcal{F}_j}
    |\langle f_j | U | f_i\rangle|^2\,d\mu(f_i)\,d\mu(f_j),
\end{equation*}
where $\mu$ is the $\mathrm{U}(n)$-invariant measure on the fibers. When
the fibers are isomorphic as $\mathrm{U}(n)$-spaces and $\mu$ is the
Haar measure, the double integral reduces to $|U_{ij}|^2$ by the
Schur orthogonality relations, establishing unistochasticity.
\end{proof}

Phase-liftability holds when the RSVP dynamics is sufficiently regular:
specifically, when the torsional components $\theta\widehat{T}$ generate
a well-defined $\mathrm{U}(n)$ action on the coarse-grained state space
and the coarse-graining $\Pi$ respects this action in the sense that the
fiber of $\Pi$ over each macrostate $C_i$ is a $\mathrm{U}(n)$-homogeneous
space. The emergence of a unistochastic transition matrix from a deeper
dynamical substrate parallels Barandes's reconstruction of quantum theory
from stochastic processes admitting a unitary lift
\cite{barandes2023,barandes2024}. Under these conditions, the path integral over
$\gamma : C_i \to C_j$ factors into a modulus-squared of a sum of
complex amplitudes, which is exactly the condition for~(\ref{eq:unistochastic}).

The quantum-mechanical Born rule is recovered when $n$ equals the dimension
of the Hilbert space and $U$ is the $S$-matrix restricted to the relevant
sector.

\subsubsection*{The Interpretive Synthesis}

The unistochastic matrix is therefore the observable shadow of hidden
rotation-scaling transport in the RSVP field space, just as the squared
modulus of a complex amplitude is the observable shadow of an unobserved
phase rotation. This synthesis unifies several seemingly disparate elements:

Needham's complex geometry shows that $i^2 = -1$ is not an axiom but a
geometric fact about quarter-turns; the imaginary unit is unsquared rotation.

Barandes's unistochastic framework shows that quantum transition probabilities
need not be postulated but can emerge from any dynamics whose coarse-graining
yields a phase-liftable stochastic kernel.

The RSVP field operator $Z = \Phi e^{\widehat{L}}$ provides a concrete
field-theoretic realization in which the unsquared object is a compound of
scaling and transport, and the Born rule emerges from discarding the
transport component under coarse-graining.

The Krein-space ghost states are a further extension of the same idea. A
ghost state in $\mathcal{K}_-$ is a direction in $M_{\mathrm{adm}}$ along
which the internal geometry has a reflection---a direction whose squared
norm is negative because the rotation it encodes is hyperbolic rather than
elliptic. When the ghost-parity symmetry ensures that the hyperbolic
contributions cancel in the trace $\mathrm{tr}(A^\dagger A)$, the ghost is
squared away along with the phases, and the observable probability is
non-negative.

The pattern is consistent across all cases: amplitudes, phases, negative
norms, and ghost sectors are all internal geometry in $M_{\mathrm{adm}}$.
Probabilities are what survives after the projection $\Pi$ discards this
geometry---after the squaring, the modulus-taking, the trace. The present
framework shares with Barandes's proposal the central structural claim that
quantum probabilities need not be fundamental: they arise from projection of
a richer underlying dynamics onto an observable transition structure
\cite{barandes2023,barandes2024}. The two frameworks differ in ontology---
Barandes grounds the derivation in classical stochastic processes, while the
admissibility framework grounds it in field-theoretic admissibility flow---
but both conclude that the Born rule is a theorem about coarse-graining
rather than an axiom about measurement. The Born rule is not a fundamental
postulate about the relationship between wave functions and measurement
outcomes; it is a theorem about what coarse-graining does to
rotation-scaling transport in a space that is generically larger, and
structurally richer, than a Hilbert space.

\subsection{Strings and the Assumption of Hilbert-Space Quantum Mechanics}

Turok draws a sharp conclusion from the Krein-space proposal: the entire
string theory programme rested on assumptions that may be unnecessarily
strong. He identifies two key assumptions behind the conclusion that quantum
gravity requires extra dimensions: first, that only theories with at most
two derivatives in the action are allowed; second, and more fundamentally,
that ``the theory lives in a Hilbert space, which means that the norms of
all quantum states are positive.''

String theory satisfies the Hilbert Assumption by construction. Its
consistency conditions---Weyl anomaly cancellation, modular invariance,
the absence of negative-norm physical states---are all requirements that
the theory project cleanly onto a positive-definite Hilbert space. The
dimension of spacetime (10 or 26) and the existence of supersymmetry are
consequences of these requirements. If the Hilbert Assumption is relaxed,
the entire edifice of constraints changes.

Turok's conclusion is direct: ``Now that we've seen that you don't need
that assumption, the whole thing has no basis.'' This may be too strong;
string theory may be correct for independent reasons. But the existence
of a UV-complete four-dimensional quantum gravity theory without strings
or extra dimensions---even in a restricted sector---removes the
\emph{necessity} argument for strings. The admissibility framework
generalizes this: the string landscape and the Krein-space quadratic
gravity theory are two distinct points in the admissibility geometry of
quantum theories that project to compatible observable physics at low
energies.


%======================================================================
\section{Quadratic Gravity as Effective Admissibility-Flow Geometry}
%======================================================================

\subsection{Geometry as the Shadow of Obstruction}

In the admissibility framework, spacetime curvature is not a fundamental
field. It arises from the obstruction to globally consistent projection:
the shadow cast by the failure of $\Pi$ to be globally flat.

Consider parallel transport of an admissibility constraint around a
closed loop in $M_{\mathrm{adm}}$. If the constraint closes upon
return, the loop contributes no curvature to the projection geometry.
If it fails to close, the discrepancy defines a curvature tensor
measuring the local non-triviality of $\Pi$. The Riemann tensor is
therefore a first-order correction to admissibility flow: a measure of
how much $\Pi$ departs from a globally flat projection. The Einstein
field equations are the conditions for admissibility-flow consistency
at leading order:
\begin{equation}
  R_{\mu\nu} - \tfrac{1}{2}g_{\mu\nu}R + \Lambda g_{\mu\nu}
  = 8\pi G\,T_{\mu\nu}.
\end{equation}
These equations equate the curvature generated by matter (right side)
with the admissibility-flow correction required by the projection geometry
(left side). General relativity is the theory of the admissibility-flow
constraint at leading order.

\subsection{Quadratic Gravity as the Second-Order Term}

If the Riemann tensor measures first-order departures of $\Pi$ from
flatness, quadratic curvature terms measure second-order departures. The
full quadratic gravity action is the first two terms of a systematic
admissibility-flow expansion:
\begin{equation}
  S = \int d^4x\,\sqrt{-g}\,
    \mathcal{L}_{\mathrm{adm}}\!\left(\Pi^{(0)}, \Pi^{(1)}, \Pi^{(2)}, \ldots\right).
\end{equation}
The Einstein--Hilbert term captures the zeroth and first-order corrections;
the quadratic terms capture the second-order correction. The theory is not
an ad hoc extension of general relativity but the natural next term in a
principled expansion.

Under this interpretation, the massive spin-2 ghost is a pole in the
propagator on $M_{\mathrm{adm}}$ arising from the second-order correction.
Whether it induces an observable pathology depends on whether $\Pi$ maps
its contribution to $M_{\mathrm{obs}}$ or to the fiber of $\Pi$. The
Turok--Bateman proposal is, in essence, a claim that---at least in the
scalar-mode limit---the ghost pole lies in the fiber of $\Pi$.

\subsection{The Effective Interface Hierarchy}

Different approximations to the admissibility-flow dynamics generate
different effective theories:

At zeroth order, the admissibility flow is trivial: spacetime is flat
and matter propagates on it without back-reaction. This is quantum field
theory in a fixed background.

At first order, the back-reaction of matter on geometry is included:
this is general relativity, or its quantization in the semiclassical
approximation.

At second order, the higher-derivative corrections are included: this is
quadratic gravity, with its ghost pole and its potential for UV completeness.

Each level is an effective interface theory---a description of $M_{\mathrm{obs}}$
obtained by truncating the admissibility-flow expansion at a given order.
The ghost is not a pathology of nature but a feature of the second-order
interface that is absent at lower-order interfaces.

\begin{principle}[Effective Interface]
Einstein gravity, quadratic gravity, and their gauge-theoretic analogues
are effective descriptions of $\Pi$ at successively higher orders in the
admissibility-flow expansion: interface approximations rather than
competing fundamental theories.
\end{principle}

\subsection{Gauge Symmetry as Fiber Structure}

Two elements $h_1, h_2 \in M_{\mathrm{adm}}$ satisfying $\Pi(h_1) = \Pi(h_2)$
represent the same observable history by different internal paths. Gauge
transformations generate exactly such equivalences: a gauge orbit in
$M_{\mathrm{adm}}$ is a set of histories projecting to a single element of
$M_{\mathrm{obs}}$. Gauge symmetry is therefore not a feature of $M_{\mathrm{obs}}$
but of the fiber of $\Pi$---the representational redundancy of the internal
description.

Diffeomorphism invariance, $\mathrm{U}(1)$ invariance in electrodynamics,
and $\mathrm{SU}(3)$ invariance in QCD are three instances of the same
phenomenon: the observable sector is carved out of a larger admissibility
manifold by a projection whose fibers are gauge orbits. The ghost-parity
symmetry of Turok and Bateman is another instance: a discrete symmetry of
$M_{\mathrm{adm}}$ needed to make the projection to $M_{\mathrm{obs}}$ well-defined.


%======================================================================
\section{Observational--Interventional Separation and Latent Fundamentalism}
%======================================================================

\subsection{The Underdetermination Theorem}

\begin{definition}[Observational Equivalence of Admissibility Manifolds]
Two admissibility manifolds $M_{\mathrm{adm}}^{(1)}$ and $M_{\mathrm{adm}}^{(2)}$
equipped with projections $\Pi_1$ and $\Pi_2$ are \emph{observationally
equivalent} if
\begin{equation}
  \Pi_1\!\bigl(M_{\mathrm{adm}}^{(1)}\bigr)
  = \Pi_2\!\bigl(M_{\mathrm{adm}}^{(2)}\bigr)
  \quad \text{as measurable spaces with probability structure.}
\end{equation}
We denote the observational equivalence class of $M_{\mathrm{adm}}$ by
$[M_{\mathrm{adm}}]_\Pi$.
\end{definition}

\begin{theorem}[Observational--Interventional Separation]
The map $[M_{\mathrm{adm}}]_\Pi \mapsto M_{\mathrm{obs}}$ is not injective:
the observational equivalence class does not determine the internal geometry
of $M_{\mathrm{adm}}$. In particular, two manifolds that differ in inner
product signature, symplectic structure, or gauge group may belong to the
same observational equivalence class.
\end{theorem}

The proof is immediate: $\Pi$ is a surjection from a larger space to a
smaller one, so its fibers are non-trivial. Its consequence is substantive. The sign of the inner product
on $M_{\mathrm{adm}}$ is underdetermined by any finite collection of
observable transition probabilities. Requiring positive-definiteness of
the inner product on $M_{\mathrm{adm}}$ is imposing a constraint on the
fiber of $\Pi$---on the unobservable interior of the theory---without
empirical justification.

\subsection{Interface Geometry via Jensen--Shannon Divergence}

The admissibility distance defined above still carries a residual commitment
to the ontology it is trying to dissolve. The inf-sup formula
\[
  d_A\!\bigl(M^{(1)}, M^{(2)}\bigr)
  = \inf_\phi \sup_h \bigl\|\Pi_1(h) - \Pi_2(\phi(h))\bigr\|
\]
compares images of the two projections, but the infimum runs over maps
$\phi$ between the underlying admissibility manifolds. It still reaches
into $M_{\mathrm{adm}}$ to do its work. A fully interface-native distance
would live entirely within $M_{\mathrm{obs}}$, requiring no reference to
the internal structure of the theories being compared.

The Jensen--Shannon divergence provides exactly this. Given two probability
distributions $P$ and $Q$ on $M_{\mathrm{obs}}$---the observable interfaces
of two admissibility manifolds, $P = \Pi_1(\mu_1)$ and $Q = \Pi_2(\mu_2)$
for some measures $\mu_i$ on $M_{\mathrm{adm}}^{(i)}$---define
\begin{equation}
  \operatorname{JSD}(P\|Q)
  = H\!\left(\frac{P+Q}{2}\right)
    - \frac{1}{2}\bigl(H(P) + H(Q)\bigr),
  \label{eq:jsd}
\end{equation}
where $H(\cdot)$ denotes the Shannon entropy. The quantity
$\operatorname{JSD}^{1/2}(P\|Q)$ is a proper metric on the space of
probability distributions \cite{endres2003}, and the entire expression is
defined without reference to either $M_{\mathrm{adm}}^{(1)}$ or
$M_{\mathrm{adm}}^{(2)}$: it depends only on the observable distributions
$P$ and $Q$.

\begin{definition}[Interface Distance]
Let $P = \Pi_1(\mu_1)$ and $Q = \Pi_2(\mu_2)$ be the observable interfaces
of two admissibility manifolds equipped with reference measures
$\mu_1, \mu_2$. The \emph{interface distance} between the two theories is
\begin{equation}
  d_{\mathrm{int}}(P, Q)
  = \operatorname{JSD}^{1/2}(P\|Q)
  = \left[
      H\!\left(\frac{P+Q}{2}\right)
      - \frac{1}{2}\bigl(H(P) + H(Q)\bigr)
    \right]^{1/2}.
  \label{eq:interface_distance}
\end{equation}
\end{definition}

The shift from $d_A$ to $d_{\mathrm{int}}$ is not merely notational. It
represents a change in what the theory of theories is \emph{about}. The
inf-sup formulation of $d_A$ still asks how close two internal structures
can be made to look; it frames indistinguishability as an obstruction to
reconstruction. The JSD formulation asks only how distinguishable two
interfaces are from within the observable domain. No hidden state. No
latent manifold. No appeal to a true underlying description. The geometry
of comparison is intrinsic to the interface itself.

The entropy form~(\ref{eq:jsd}) admits a decomposition that connects
directly to the admissibility program:
\begin{equation}
  \mathcal{D}(P, Q)
  \;=\;
  H\!\left(\frac{P+Q}{2}\right)
  - \frac{1}{2}\bigl(H(P) + H(Q)\bigr)
  \;=\;
  \operatorname{JSD}(P\|Q).
\end{equation}
The first term measures the entropy of the \emph{merged} interface---the
distribution seen by an observer who cannot distinguish which theory
generated a given outcome. The second term measures the average entropy of
the \emph{separate} interfaces. Their difference is exactly the information
that disappears when the distinction between the two interfaces is forgotten.
We may call $\mathcal{D}(P,Q)$ the \emph{distinction functional}: it
measures how much observational content is carried by the difference between
$P$ and $Q$. The Jensen--Shannon divergence is its metric realisation.

\subsubsection*{Gauge Orbits as JSD-Zero Equivalence Classes}

The JSD formulation provides a purely interface-native characterisation of
the structures previously described in terms of internal geometry. Two
histories $h, h' \in M_{\mathrm{adm}}$ are \emph{observationally
indistinguishable} when
\begin{equation}
  \operatorname{JSD}\bigl(\Pi(h),\, \Pi(h')\bigr) = 0.
\end{equation}
This is the correct definition of a gauge orbit from the interface
perspective: a gauge transformation is a deformation of the internal
history that induces zero change in the distinction functional. The
internal structure of the gauge orbit---whether it is generated by a Lie
group, a BRST operator, a ghost-parity discrete symmetry, or the harmonic
gauge symmetry of the dimension-zero scalars---is irrelevant to its
observable characterisation. All that matters is that the deformation
annihilates $\mathcal{D}$.

Likewise, a \emph{hidden admissibility direction} in the sense of
Definition~3.1 can now be given a metric characterisation: it is a
deformation of the internal theory along which the interface distance
remains arbitrarily small,
\begin{equation}
  \lim_{\epsilon \to 0}
  d_{\mathrm{int}}\!\bigl(\Pi(h),\, \Pi(h + \epsilon\, \delta h)\bigr)
  = 0
  \quad \text{while} \quad \delta h \neq 0.
\end{equation}
The ghost sector of a Krein-space theory is a hidden admissibility direction
in exactly this sense: the deformation of the internal state space that
adds a $\mathcal{K}_-$ component does not change the interface distance
between observable transition probabilities, because ghost-parity symmetry
ensures the additional contribution cancels in the traced probability
functional.

\subsubsection*{Fisher Information Geometry of the Interface Space}

For a parametric family of projections $\Pi_\lambda$, the interface
distance induces a Riemannian structure on the parameter space. For
infinitesimally separated interfaces $\Pi_\lambda$ and $\Pi_{\lambda +
\delta\lambda}$, the squared interface distance expands as
\begin{equation}
  d_{\mathrm{int}}^2\!\bigl(\Pi_\lambda,\, \Pi_{\lambda+\delta\lambda}\bigr)
  = \frac{1}{8}\,g_{ij}(\lambda)\,\delta\lambda^i\,\delta\lambda^j
    + O(\delta\lambda^3),
\end{equation}
where the metric tensor is
\begin{equation}
  g_{ij}(\lambda)
  = \left.
      \frac{\partial^2}{\partial\lambda^i\,\partial\lambda^j}
      \operatorname{JSD}^2\!\bigl(\Pi_\lambda,\, \Pi_{\lambda+\delta\lambda}\bigr)
    \right|_{\delta\lambda=0}.
  \label{eq:fisher_metric}
\end{equation}
For small separations this reduces to the Fisher information metric on the
manifold of probability distributions \cite{amari2000}, establishing that
$(\mathcal{M}, d_{\mathrm{int}})$ is locally a statistical manifold in the
sense of information geometry.

This is not the curvature of spacetime. It is not the curvature of a
quantum state manifold. It is the curvature of \emph{distinguishability}:
the rate at which the observable interfaces of nearby theories diverge as a
function of the theory parameters. The admissibility geometry of quantum
theories---the space $(\mathcal{M}, d_{\mathrm{int}})$---is therefore a
well-defined Riemannian manifold whose geodesics are paths through theory
space that minimise the total accumulated distinguishability between
successive observable interfaces.

\subsubsection*{Loss Becomes Relational}

The philosophical consequence of adopting $d_{\mathrm{int}}$ as the
primitive metric on theory space is a reorientation that runs deeper than
any particular result about ghosts or quadratic gravity.

Under the inf-sup formulation of $d_A$, the question one asks about two
theories is: \emph{what hidden content is inaccessible?} The comparison
reaches into the fibers of $\Pi$ and asks how different the internal
structures are from each other. Loss is ontological: something exists in
$M_{\mathrm{adm}}$ that does not appear in $M_{\mathrm{obs}}$.

Under the JSD formulation of $d_{\mathrm{int}}$, the question becomes:
\emph{how distinguishable are the interfaces?} The comparison lives entirely
within $M_{\mathrm{obs}}$ and asks how much observational content separates
the two probability distributions. Loss is relational: the distinction
functional $\mathcal{D}(P,Q)$ measures how much information disappears when
one stops distinguishing between the two interfaces. There is no commitment
to a hidden object that is being lost---only to a geometry of distinction
that measures how much the forgetting costs.

This reorientation places the framework in a precise relationship to the
admissibility program as a whole. The program began with the conviction
that reachability, admissibility, and constraint closure are the correct
ontological primitives. The JSD formulation now suggests a further
reduction: the primitive is not the admissibility manifold $M_{\mathrm{adm}}$
but the distinguishability structure on $M_{\mathrm{obs}}$. The admissibility
manifold is a representational tool for generating and organising interface
distributions; the geometry of those distributions, measured by
$d_{\mathrm{int}}$, is what physics is ultimately about. In this sense the
framework converges on an information-geometric version of the distinction
program:
\begin{equation}
  \text{Reality}
  \;\longrightarrow\;
  \text{Distinctions}
  \;\longrightarrow\;
  \text{Distances between distinctions,}
\end{equation}
where the final object is the metric space $(\mathcal{M}, d_{\mathrm{int}})$
and the Jensen--Shannon divergence is the first mathematically natural
metric that lives entirely at that level without requiring access to
anything underneath.

\begin{remark}
Theories in the same observational equivalence class (Definition~8.1) have
$d_{\mathrm{int}} = 0$ by construction. Theories at $d_{\mathrm{int}} > 0$
make empirically distinguishable predictions. The Derived Interface
Conjecture of Section~9 asserts that quadratic gravity and Hilbert-space
quantum gravity satisfy $d_{\mathrm{int}} = 0$ below the Planck scale, and
$d_{\mathrm{int}} > 0$ above it.
\end{remark}

\subsection{Against Latent Fundamentalism}

\emph{Latent fundamentalism} is the tendency to grant ontological status
to internal coordinates merely because they appear in a successful model.
It manifests systematically in the physics of higher-derivative theories.

In ghost arguments: the internal coordinates of the extended phase space
are granted direct physical interpretation, and their unbounded range is
treated as a certificate of physical instability rather than as a coordinate
property of the representation.

In the Hilbert-space requirement: the positive-definiteness of the inner
product on the state space is treated as a necessary condition for
physical admissibility when it is at most a sufficient one.

In multiverse arguments: the proliferation of branches in the quantum
state of the universe is treated as physically real on the ground that the
wave function is the fundamental object of quantum mechanics. But the wave
function is an element of $M_{\mathrm{adm}}$, not of $M_{\mathrm{obs}}$.

In each case the admissibility perspective offers the same corrective:
internal coordinates are instruments of representation. Their pathologies
are coordinate pathologies until proven otherwise. The physical question
is always: what does $\Pi$ map them to?

Turok makes this argument directly, and extends it to the assumptions behind
string theory: the requirement that theories live in Hilbert spaces was
``probably the strongest assumption'' behind the conclusion that quantum
gravity requires extra dimensions. Now that an example exists of a
UV-complete quantum theory of gravity that does not live in a Hilbert
space, the necessity argument for strings loses its basis. The admissibility
framework generalizes this critique: the question of which internal structure
corresponds to the observable world is empirical, not a priori.

\subsection{A Systematic Table of Coordinate and Observable Pathologies}

The distinction between coordinate and observable pathologies applies
across a range of well-understood physical contexts. The following table
illustrates the pattern:

\medskip
{\small
\begin{center}
\begin{tabular}{p{3.2cm}p{4.2cm}p{4.2cm}}
\toprule
Context & Coordinate Pathology & Status in $M_{\mathrm{obs}}$ \\
\midrule
Lorenz-gauge QED & Negative-norm photons & Absent from physical states \\
BRST gauge theories & Off-shell Faddeev--Popov ghosts & Absent from BRST cohomology \\
Pais--Uhlenbeck model & Unbounded Hamiltonian & Bounded periodic trajectories \\
Ostrogradsky (gravity) & Runaway phase-space direction & Cosmological expansion \\
Krein-space QFT & Indefinite inner product & Positive traced probabilities \\
Spin-2 ghost (QG) & Negative propagator pole & Status under investigation \\
Gauge orbits & Representational redundancy & Single observable history \\
Wave function branches & Proliferating $M_{\mathrm{adm}}$ directions & Observable record underdetermined \\
\bottomrule
\end{tabular}
\end{center}}
\medskip

The table illustrates that coordinate pathologies are the rule rather than
the exception in well-understood physical theories. The history of physics
is substantially a history of learning to identify which features of a
formalism are representational and which are physical.

\subsection{The Philosophical Position of Quadratic Gravity}

The gravitational entropy program of Turok connects to the admissibility
perspective in a striking way. Rather than deriving the large-scale
homogeneity and isotropy of the universe from special initial conditions
or a smoothing mechanism like inflation, Turok asks whether the observed
universe is simply a \emph{typical} element of $M_{\mathrm{obs}}$ for the
relevant $M_{\mathrm{adm}}$.

This is exactly the strategy of the admissibility program applied to
cosmology: rather than seeking a dynamical mechanism that drives the
universe toward observed states, one asks how large the admissible
trajectory space is for states consistent with observation. In the
admissibility framework, the gravitational entropy associated with a
cosmological history is the logarithm of the volume of the corresponding
admissibility set:
\begin{equation}
  S_{\mathrm{grav}} \sim \log\,\mathrm{Vol}(\mathcal{A}_H).
\end{equation}
The observed universe is typical if $\mathcal{A}_H$ is large---if many
distinct elements of $M_{\mathrm{adm}}$ project to the observed large-scale
structure. The problem of cosmological initial conditions is thereby
reframed not as a question about dynamics but as a question about the
geometry of $\Pi$.


%======================================================================
\section{The Derived Interface Conjecture}
%======================================================================

\subsection{Statement}

The arguments of the preceding sections converge on a single conjecture.
We state it in general and then in the specific form relevant to quantum
gravity.

\begin{conjecture}[Derived Interface, General Form]
Observable physical theories are effective interfaces generated by
projection from a larger admissibility manifold $M_{\mathrm{adm}}$.
The structural features of such theories---Hilbert space, positive inner
product, gauge symmetry, metric positivity---are properties of the
interface, not of $M_{\mathrm{adm}}$ itself.
\end{conjecture}

\begin{conjecture}[Derived Interface, Quantum Gravity Form]
Observable spacetime theories and their quantizations are effective
interfaces of the following kind:
\begin{enumerate}[label=(\roman*)]
  \item Einstein gravity is the leading-order interface description of
        admissibility-flow consistency on $M_{\mathrm{adm}}$.
  \item Quadratic gravity is the second-order description, retaining
        first admissibility-curvature corrections. The massive ghost pole
        corresponds to a negative-metric direction in the fiber of $\Pi$
        that does not reach $M_{\mathrm{obs}}$ at sub-Planckian energies.
  \item Krein-space sectors correspond to hidden admissibility directions
        in $M_{\mathrm{adm}}$ that survive internally but cancel under $\Pi$
        at accessible energies, provided ghost-parity symmetry is preserved.
  \item Gauge symmetry encodes the fiber structure of $\Pi$: redundant
        directions in $M_{\mathrm{adm}}$ mapping to a single element of
        $M_{\mathrm{obs}}$.
  \item The Hilbert Assumption is the zeroth-order interface approximation,
        valid when the fiber of $\Pi$ is negligible and all internal degrees
        of freedom are observable.
\end{enumerate}
\end{conjecture}

\subsection{Falsifiability and Empirical Content}

The Derived Interface Conjecture is falsifiable. A physical
effect---a scattering amplitude, a cosmological observable, a precision
measurement---that could not be reproduced by any admissibility-consistent
projection from any $M_{\mathrm{adm}}$ would refute the general form.

For the quantum gravity form, the most direct test is the one Turok
identifies: whether the ghost-parity symmetry extends to the spin-2 sector
of quadratic gravity. If it does, the spin-2 ghost is a hidden admissibility
direction and the specific Conjecture~9.2(ii) gains strong support. If
it does not, the spin-2 ghost generates observable negative probabilities
at energies where it goes on-shell, and the theory as stated must be
modified.

The constructive reading of the conjecture defines a research programme:
given $M_{\mathrm{obs}}$ corresponding to current observations, what is
the space of admissibility manifolds $M_{\mathrm{adm}}$ that project onto
it? This \emph{admissibility geometry of quantum theories} has Hilbert-space
quantum mechanics as one distinguished point. Krein-space theories are
neighboring points differing in the inner-product signature. Higher-derivative
theories are points at greater distance, differing in kinetic structure.
Exploring the geometry of this space---which regions project onto compatible
$M_{\mathrm{obs}}$, which do not, what distinguishes neighboring points
empirically---is the task that the admissibility framework makes precise.

\subsection{On the Necessity of Positive Norm}

The compact summary of the paper's principal conclusion is:
\begin{equation}
  \text{positive norm on } M_{\mathrm{adm}}
  \;\Longrightarrow\;
  \text{observable admissibility,}
\end{equation}
but the converse fails in general. Positive norm is a sufficient condition,
not a necessary one. Hilbert-space quantum mechanics occupies a privileged
and well-understood region of the admissibility geometry, but it is a
region, not the entire geometry. Theories with indefinite inner product on
$M_{\mathrm{adm}}$ may be perfectly admissible in $M_{\mathrm{obs}}$, and
the burden of proof lies with those who claim otherwise.

The implication for the programme of quantizing gravity is substantial. The
conclusion that gravity requires extra dimensions or strings was predicated
on the assumption that admissibility is imposed on $M_{\mathrm{adm}}$ rather
than on $M_{\mathrm{obs}}$. Relaxing this assumption does not guarantee that
quadratic gravity is the correct theory of quantum gravity; it may not be.
But it removes the a priori impossibility of four-dimensional quantum gravity
without exotic structure, and restores the question to the domain of
empirical investigation where it belongs.


%======================================================================
\section{Conclusion}
%======================================================================

We have argued that the standard dismissal of higher-derivative quantum
gravity rests on the Hilbert Assumption: the identification of the full
internal state space of a theory with its observable sector. Relaxing
this assumption opens a coherent framework---the admissibility manifold
hierarchy---in which negative-norm sectors, Krein-space structures, and
quadratic gravity appear not as pathologies but as coordinate features of
a larger internal space projecting onto the observable world.

The central structural results are:

The \emph{Projection Sufficiency Principle}: physical constraints are
conditions on $M_{\mathrm{obs}}$, not on $M_{\mathrm{adm}}$.

The \emph{Projection Pathology Separation Theorem}: coordinate pathologies
in $M_{\mathrm{adm}}$: negative inner products, unbounded
Hamiltonians, gauge redundancies---need not propagate to observable
pathologies in $M_{\mathrm{obs}}$.

The \emph{Observable Admissibility Theorem}: a Krein-space quantum theory
with ghost-parity symmetry produces non-negative observable probabilities
summing to unity, regardless of the sign of individual ghost-state inner
products. This is the content of the Turok--Bateman proposal, derived from
the general framework.

The \emph{Observable Stability Principle}: physical stability is determined
by the admissibility of projected trajectories, not by the sign of the
Hamiltonian. The Ostrogradsky theorem is a statement about $M_{\mathrm{adm}}$;
its implications for $M_{\mathrm{obs}}$ must be determined via $\Pi$.

The \emph{Observational--Interventional Separation Theorem}: the observable
record underdetermines the structure of $M_{\mathrm{adm}}$, establishing a
principle of ontological restraint on the interpretation of internal
coordinates.

These results converge on the \emph{Derived Interface Conjecture}: observable
spacetime theories are effective interfaces generated by projection from a
larger admissibility manifold. Einstein gravity and quadratic gravity are
the first two terms of a systematic admissibility-flow expansion. Ghost
sectors are hidden admissibility directions in the fiber of $\Pi$. Gauge
symmetry is fiber structure. The Hilbert Assumption is the zeroth-order
interface approximation.

Hilbert-space quantum mechanics remains the best-confirmed physical
framework we possess. The argument here is not that it is wrong but that
it may occupy a privileged region inside a larger admissibility geometry
rather than constituting that geometry in its entirety. Krein-space
theories, higher-derivative theories, and quadratic quantum gravity represent
neighboring admissible sectors of that geometry---sectors that the observable
interface has not, at accessible energies, distinguished from the interior.

The ghost need not be exorcised. It need only be projected.

\bigskip

%----------------------------------------------------------------------

%======================================================================
\section*{Appendix A: The Dimension-Zero Scalar Programme}
\addcontentsline{toc}{section}{Appendix A: The Dimension-Zero Scalar Programme}
%======================================================================

The arguments of this paper connect to a research programme, developed by Boyle and Turok, that
addresses several foundational puzzles simultaneously through the introduction of \emph{dimension-zero
scalar fields}. We outline this programme here because it exemplifies, at a concrete level, how the
admissibility framework leads to new physics rather than mere reinterpretation of existing results.

\subsection*{A.1 Ordinary versus Dimension-Zero Scalars}

An ordinary conformally coupled scalar $\varphi$ of mass dimension one is described by the action
\begin{equation}
  S_2[\varphi] = \tfrac{1}{2}\int d^4x\,\sqrt{g}\,\varphi\,\Delta_2\,\varphi,
\end{equation}
where $\Delta_2 = -\Box + \tfrac{1}{6}R$ is the unique conformally invariant second-order
operator. The action is invariant under $g_{\mu\nu} \to \Omega^2 g_{\mu\nu}$,
$\varphi \to \Omega^{-1}\varphi$.

A dimension-zero conformally coupled scalar is instead described by
\begin{equation}
  S_4[\varphi] = \pm\tfrac{1}{2}\int d^4x\,\sqrt{\pm g}\,\varphi\,\Delta_4\,\varphi,
\end{equation}
where the sign convention follows from demanding positivity of the Euclidean action and then
Wick-rotating, and $\Delta_4$ is the unique conformally invariant fourth-order operator,
\begin{equation}
  \Delta_4 = \Box^2 + 2R^{\mu\nu}\nabla_\mu\nabla_\nu
             - \tfrac{2}{3}R\,\Box + \tfrac{1}{3}(\nabla^\mu R)\nabla_\mu.
\end{equation}
This action is invariant under $g_{\mu\nu} \to \Omega^2 g_{\mu\nu}$,
$\varphi \to \Omega^0\varphi$: the field does not transform under Weyl rescalings,
which accounts for the dimension-zero designation.

In flat spacetime the action reduces to
\begin{equation}
  S_4 = \pm\tfrac{1}{2}\int d^4x\,\varphi\,\Box^2\varphi
       = \pm\tfrac{1}{2}\int d^4x\,(\Box\varphi)^2.
\end{equation}
This is a fourth-order theory, so one should expect ghost states in the propagator. However,
as Bogoliubov et al.\ established and Boyle and Turok exploited, the action~(A.3) possesses
an infinite-dimensional gauge symmetry
\begin{equation}
  \varphi \;\longrightarrow\; \varphi + \chi, \quad \Box\chi = 0.
\end{equation}
This symmetry shifts $\varphi$ by any harmonic function, and once it is properly taken
into account, the theory has \emph{no local degrees of freedom at all}. Every local gauge-invariant
operator commutes, and the theory possesses a unique physical state: the vacuum. This is the
decisive fact about dimension-zero scalars. Their ghost problem is not resolved by projecting onto
a positive Hilbert subspace but by the gauge symmetry making all local degrees of freedom
pure gauge.

\subsection*{A.2 The Vacuum Energy and the Weyl Anomaly}

The dimension-zero scalar contributes to the vacuum energy per Fourier mode with a
coefficient of opposite sign to ordinary scalars. Combining the contributions from
dimension-one scalars ($n_0$), Weyl fermions ($n_{1/2}$), gauge bosons ($n_1$), and
dimension-zero scalars ($n_0'$), one finds
\begin{equation}
  E_k^{(0)} = \frac{\hbar k}{2}\bigl[n_0 - 2n_{1/2} + 2n_1 + 2n_0'\bigr].
\end{equation}
Similarly, the trace of the stress-energy tensor in curved spacetime---the Weyl anomaly---takes
the form
\begin{equation}
  \langle T^\mu_\mu \rangle = c\,C^2 - a\,E + \xi R,
\end{equation}
where $C^2 = C_{\mu\nu\rho\sigma}C^{\mu\nu\rho\sigma}$ is the square of the Weyl tensor,
$E = R_{\mu\nu\rho\sigma}R^{\mu\nu\rho\sigma} - 4R_{\mu\nu}R^{\mu\nu} + R^2$ is the Euler
density, and the coefficients at one loop are
\begin{align}
  a &= \frac{1}{360(4\pi)^2}\bigl[n_0 + \tfrac{11}{2}n_{1/2} + 62n_1 - 28n_0'\bigr],\\
  c &= \frac{1}{120(4\pi)^2}\bigl[n_0 + 3n_{1/2} + 12n_1 - 8n_0'\bigr].
\end{align}
The key observation of Boyle and Turok is that dimension-zero scalars contribute
\emph{negatively} to both $a$ and $c$, providing an alternative to supersymmetry as a
mechanism for Weyl anomaly cancellation.

\subsection*{A.3 The Three-Generation Prediction}

Consider the standard model gauge group $\mathrm{SU}(3)\times\mathrm{SU}(2)\times\mathrm{U}(1)$
with $n_1 = 12$ gauge bosons, and suppose the Higgs is composite or emergent so that $n_0 = 0$.
The conditions for simultaneous cancellation of the vacuum energy and both terms in the
Weyl anomaly---i.e.\ setting $E_k^{(0)} = 0$, $a = 0$, $c = 0$---constitute three equations
in the four unknowns $\{n_0, n_{1/2}, n_1, n_0'\}$. The unique solution is
\begin{equation}
  n_{1/2} = 4n_1, \qquad n_0' = 3n_1, \qquad n_0 = 0.
\end{equation}
For $n_1 = 12$, this gives $n_{1/2} = 48 = 3 \times 16$: exactly three generations of standard
model fermions, each including a right-handed neutrino. The number of dimension-zero scalars
required is $n_0' = 36$.

This is a remarkable result: the number of fermion generations is not a free parameter of the
model but a consequence of requiring simultaneous cancellation of three independent gravitational
anomalies. The result is non-trivial in that $E_k^{(0)}$, $a$, and $c$ are completely independent
linear combinations of $\{n_0, n_{1/2}, n_1, n_0'\}$; their simultaneous vanishing could
easily have been inconsistent.

\subsection*{A.4 Admissibility Interpretation}

In the language of this paper, the dimension-zero scalar programme is exactly an instance of
the Projection Sufficiency Principle. The fourth-order action in $M_{\mathrm{adm}}$ carries
coordinate pathologies---ghost poles in the propagator, an indefinite inner product. But once
the gauge symmetry $\varphi \to \varphi + \chi$ is accounted for, the projection $\Pi$ to
$M_{\mathrm{obs}}$ eliminates all local degrees of freedom. The observable sector is simply the
vacuum, with its scale-invariant fluctuation spectrum. The ghost never reaches $M_{\mathrm{obs}}$.

The debate between Boyle--Turok and their critics (Cline and Hell, 2025) concerns exactly
whether the gauge symmetry is sufficient to remove the ghost from $M_{\mathrm{obs}}$, or whether
the ghost survives under certain background conditions (curved spacetime, time-dependent background
field). In the admissibility framework, this is the question of whether the projection $\Pi$ is
sufficiently robust across all physical backgrounds, not merely in flat spacetime. It is a concrete
instance of the general question posed in Conjecture~9.2: does the relevant symmetry---whether
ghost-parity or the harmonic-function gauge symmetry---extend from the restricted sector to the
full theory?


%======================================================================
\section*{Appendix B: The CPT-Symmetric Universe and Gravitational Entropy}
\addcontentsline{toc}{section}{Appendix B: The CPT-Symmetric Universe and Gravitational Entropy}
%======================================================================

\subsection*{B.1 The CPT Universe}

Boyle, Finn, and Turok proposed in 2018 that the state of the universe does not spontaneously
violate $CPT$: the universe after the Big Bang is the $CPT$ image of the universe before it,
both classically and quantum mechanically. The pre- and post-bang epochs comprise a
universe/anti-universe pair emerging from a conformal singularity at $\tau = 0$.

The key classical observation is that the radiation-dominated flat FRW metric takes the form
$ds^2 = a^2(\tau)(-d\tau^2 + d\mathbf{x}^2)$ with $a(\tau) \propto \tau$ near the bang.
Analytically continuing across $\tau = 0$ reveals a new isometry: time reversal $\tau \to -\tau$.
The $CPT$ symmetry condition on the tetrad is
\begin{equation}
  e^a_\mu(\tau, \mathbf{x}) = -e^a_\mu(-\tau, \mathbf{x}),
\end{equation}
which immediately implies that $a(\tau)$ is an odd function with a simple zero at $\tau = 0$.

The quantum consequences are substantial. In a general curved spacetime, the choice of vacuum is
ambiguous: different observers at different epochs define inequivalent vacua related by Bogoliubov
transformations. The $CPT$ symmetry, however, selects a unique preferred vacuum $|0_0\rangle$
invariant under $CPT|0_0\rangle = |0_0\rangle$. This $CPT$-symmetric vacuum is the geometric
mean of the vacua preferred by observers in the far past and far future.

The $CPT$ condition on the scalar perturbations selects the regular, growing mode and eliminates
the singular decaying mode near $\tau = 0$---precisely the boundary condition responsible for
the acoustic oscillations in the CMB power spectrum that are usually attributed to inflation.
For tensor perturbations, the same condition eliminates primordial gravitational waves: the
CPT-symmetric universe predicts $r = 0$ to this order, distinguishing it from inflationary models.

\subsection*{B.2 The Gravitational Entropy Calculation}

Boyle and Turok developed a thermodynamic explanation for the observed flatness, homogeneity,
and isotropy of the universe based on a new calculation of gravitational entropy. The central
observation is that cosmological solutions with radiation and a cosmological constant can be
analytically continued to Euclidean instantons via a combined conformal and Wick rotation
($CW$ rotation), in which both $n \to -iN$ (standard Wick rotation) and $a(t) \to ib(t)$
(conformal rotation). The latter is required because $a(\tau)$ has a zero at the bang and
becomes imaginary under standard Wick rotation.

The gravitational entropy of a cosmological solution is the semiclassical exponent of the
Euclidean path integral:
\begin{equation}
  S_g = iS_{\mathrm{Lorentzian}} = -S_E.
\end{equation}
For a universe with cosmological constant $\rho_\Lambda$, radiation entropy $S_r$, and
spatial curvature $\kappa$, the calculation yields
\begin{equation}
  S_g \sim S_\Lambda^{1/4} S_r,
\end{equation}
where $S_\Lambda = 24\pi^2 M_{\mathrm{Pl}}^4 / \rho_\Lambda$ is the de Sitter entropy.

The physical interpretation is as follows. The partition function of the cosmological ensemble
is $Z = e^{S_g}$. The total entropy is maximized when the universe is flat, homogeneous, and
isotropic with a small positive cosmological constant. The observed large-scale geometry of the
cosmos is therefore typical rather than fine-tuned: it is the most probable macrostate of the
gravitational ensemble. Inflation is no longer needed as a dynamical smoothing mechanism.

In the language of the admissibility framework, the gravitational entropy is the logarithm of
the volume of the admissibility set of cosmological histories consistent with the observed
large-scale parameters:
\begin{equation}
  S_g \sim \log\,\mathrm{Vol}(\mathcal{A}_H^{\mathrm{cosm}}).
\end{equation}
The observed universe is typical because $\mathcal{A}_H^{\mathrm{cosm}}$ is large: many
distinct elements of $M_{\mathrm{adm}}$ project to the observed flat, homogeneous observable
sector. This is a reachability-volume argument applied to cosmological histories, and it
connects the gravitational entropy programme directly to the admissibility geometry of
Section~9.

\subsection*{B.3 Dark Matter as a CPT Consequence}

In the CPT-symmetric universe, the right-handed neutrinos (which must exist by the
three-generation prediction of Appendix~A) are produced from the CPT-symmetric vacuum
by the same mechanism that produces Hawking radiation from black holes. The $\nu_R$ field
equation is regular at the bang; the CPT-symmetric vacuum state then determines a unique
initial condition for $\nu_R$ production, from which the density of right-handed neutrinos
in the observable universe can be calculated without free parameters.

If one of the three right-handed neutrinos is stable---which follows from a $\mathbb{Z}_2$
symmetry that simultaneously makes the lightest neutrino massless---its dark matter density
matches the observed $\Omega_{\mathrm{DM}}$ when its mass is $M \approx 5 \times 10^8\,\mathrm{GeV}$.
This is a prediction, not a fit: the dark matter density is determined by the same CPT boundary
conditions that fix the CMB power spectrum.

The theory also predicts: (i) the three light neutrinos are Majorana and permit neutrinoless
double beta decay; (ii) the lightest neutrino is massless; (iii) no primordial long-wavelength
gravitational waves ($r = 0$ at this order).


%======================================================================
\section*{Appendix C: The Criticism and the Admissibility Response}
\addcontentsline{toc}{section}{Appendix C: The Criticism and the Admissibility Response}
%======================================================================

The criticism of Cline and Hell (2025) targets the dimension-zero scalar programme directly.
Their argument is that the ghost in the $S_4[\varphi]$ theory is not eliminated by the harmonic
gauge symmetry and persists as an observable instability, particularly in curved spacetime
backgrounds.

Their central technical point is the following. In flat spacetime, the Hubbard--Stratonovich
decomposition of $S_4$ introduces an auxiliary field $\sigma$ and yields the two-derivative
Lagrangian
\begin{equation}
  S \simeq \int d^4x\!\left(-\partial_\mu\sigma\,\partial^\mu\varphi - \tfrac{1}{2}\sigma^2\right),
\end{equation}
whose kinetic matrix has eigenvalues $\pm 1$, revealing a ghost. The corresponding Hamiltonian
density
\begin{equation}
  \mathcal{H} = \dot\varphi\,\dot\sigma + \nabla\varphi\cdot\nabla\sigma + \tfrac{1}{2}\sigma^2
\end{equation}
is unbounded below, and explicit wave-like solutions demonstrate locally unbounded energy density.
In FRW backgrounds, Cline and Hell show that the coefficient $c_1$ of the kinetic term for
scalar perturbations becomes negative at large wave numbers, indicating superluminal ghost modes.
They also argue that coupling to standard model fields induces a confining fifth force that rules
out the dimension-zero scalars as a source of primordial perturbations.

From the admissibility perspective, the disagreement is a disagreement about whether the projection
$\Pi$ maps the ghost to the fiber of $\Pi$ or to $M_{\mathrm{obs}}$. Boyle and Turok's claim is
that the gauge symmetry $\varphi \to \varphi + \chi$ removes the ghost from any physical prediction
because all local gauge-invariant operators commute. Cline and Hell's claim is that this argument
applies only in flat spacetime and that the gauge symmetry does not extend to FRW backgrounds in
a way that continues to eliminate the ghost from the physical spectrum.

This is precisely the kind of case-by-case question that the Projection Pathology Separation
Theorem identifies as needing separate analysis. The theorem establishes that coordinate
pathologies \emph{can} fail to propagate to $M_{\mathrm{obs}}$; it does not establish that they
always do so. Whether the harmonic gauge symmetry is robust against curved-background perturbations
is a technical question that the general framework leaves open, to be resolved by detailed
calculation.

The exchange between Boyle--Turok and Cline--Hell is therefore a concrete instantiation of the
admissibility programme in action: the ghost in $M_{\mathrm{adm}}$ is identified; the question
is whether a specific symmetry maps it to the fiber of $\Pi$ under all physically relevant
projection conditions. The current status is that this question is not fully resolved, and the
outcome will determine whether the dimension-zero scalar programme is viable.

This is characteristic of the general situation identified in Section~9: the admissibility
framework does not make theories easy; it makes the right question precise. Whether a ghost is
an observable pathology or a hidden admissibility direction is always a question about $\Pi$,
and answering it requires both the general framework and the specific technical analysis of the
theory in question.


%======================================================================
\section*{Appendix D: Scale-Invariant Primordial Perturbations Without Inflation}
\addcontentsline{toc}{section}{Appendix D: Scale-Invariant Primordial Perturbations Without Inflation}
%======================================================================

One of the most striking predictions of the combined programme is that the dimension-zero
scalars can generate the observed scale-invariant spectrum of primordial density perturbations
\emph{without inflation}. We sketch the mechanism here because it connects the UV structure
of the theory (the dimension-zero scalar action) to the IR observations (the CMB power spectrum),
illustrating how the admissibility framework connects across scales.

The two-point function of the dimension-zero scalar in flat spacetime is
\begin{equation}
  \langle\varphi(t,\mathbf{x})\,\varphi(t,\mathbf{x}')\rangle
  = \int\!\frac{d^3k}{(2\pi)^3}\,e^{i\mathbf{k}\cdot(\mathbf{x}-\mathbf{x}')}\,\frac{1}{4k^3}.
\end{equation}
The factor $1/4k^3$ is scale invariant: the power spectrum $P(k) \propto k^{-3}$ means that
each logarithmic interval of $k$ contributes equally to $\langle\varphi^2\rangle$. This is
\emph{Harrison--Zel'dovich} scale invariance, the same spectrum observed in the CMB, arising
here from the conformal invariance of $S_4$ rather than from inflationary dynamics.

In a flat FRW spacetime, conformal invariance ensures that the scale factor $a(\tau)$ does not
appear in the action~(A.3), so the same calculation applies cosmologically: the vacuum
fluctuations of $\varphi$ are cosmologically scale invariant.

The coupling between the dimension-zero scalars and standard model running couplings generates
a slight red tilt. The running coupling constant $g(\mu)$ varies with energy scale, and
this variation introduces a small breaking of scale invariance. For the standard model gauge
couplings extrapolated to the Planck scale, the predicted spectral index is approximately
\begin{equation}
  n_s \approx 0.958,
\end{equation}
in good agreement with the Planck satellite measurement $n_s = 0.959 \pm 0.006$. The tensor-to-scalar
ratio is predicted to be zero (no primordial gravitational waves), consistent with the current
upper bound $r < 0.036$ from BICEP/Keck.

The connection to the admissibility framework is via the unistochastic picture of Section~6.
The primordial scalar perturbations are transition probabilities from the initial vacuum
configuration to the late-time CMB configuration, computed by integrating over field trajectories
in $M_{\mathrm{adm}}$ weighted by $|\Phi|^2$. The scale invariance of the perturbation
spectrum reflects the conformal invariance of the $S_4$ dynamics on $M_{\mathrm{adm}}$;
the slight red tilt reflects the running of couplings that breaks conformal invariance at
finite energy. The observable CMB spectrum is $\Pi$ applied to this ensemble of field
trajectories.



%======================================================================
\section*{Appendix E: Consolidated Proof Scaffold}
\addcontentsline{toc}{section}{Appendix E: Consolidated Proof Scaffold}
%======================================================================

The following collects tightened proofs of the principal results, filling
gaps identified in the main text and adding several results that complete
the synthesis. Notational conventions follow the main text.

\subsection*{E.1 Projection Sufficiency as a Consequence of Definition}

\begin{theorem}[Projection Sufficiency, formal version]
Let $A_O : M_{\mathrm{obs}} \to \{0,1\}$ be the observable admissibility
predicate, and let $A_H(h) = A_O(\Pi(h))$ be its pullback to
$M_{\mathrm{adm}}$. If physical predictions depend only on $A_O(\Pi(h))$,
then no condition on $h \in M_{\mathrm{adm}}$ is physically necessary
unless it changes $\Pi(h)$.
\end{theorem}

\begin{proof}
Let $h, h' \in M_{\mathrm{adm}}$ satisfy $\Pi(h) = \Pi(h')$. Then
$A_H(h) = A_O(\Pi(h)) = A_O(\Pi(h')) = A_H(h')$.
Any property distinguishing $h$ from $h'$ while leaving $\Pi(h)$ fixed
has no observable effect. Physical constraints may therefore be imposed
entirely on $M_{\mathrm{obs}}$ without imposing them on every coordinate
direction of $M_{\mathrm{adm}}$.
\end{proof}

\begin{corollary}[Hilbert Assumption as injective special case]
Positive-definite Hilbert quantum mechanics is the special case in which
$\Pi$ is injective. When $\Pi$ is injective, $\Pi(h) = \Pi(h')$ implies
$h = h'$, so every fiber is a singleton and there are no hidden
admissibility directions. Every condition on internal coordinates is
immediately a condition on observable predictions.
\end{corollary}

\begin{proof}
Injectivity of $\Pi$ means $\Pi^{-1}(o) = \{h\}$ for each $o \in M_{\mathrm{obs}}$.
There is therefore no $h \neq h'$ with $\Pi(h) = \Pi(h')$; the fiber is
trivial, the internal coordinate is fully observable, and the Hilbert
Assumption holds.
\end{proof}

\subsection*{E.2 Ghost States: Negative Norm Is Not Negative Probability}

The following makes precise the claim that negative norm is a coordinate
pathology rather than an observable pathology, using the Krein space
structure of Section~3.

\begin{proposition}
Let $g \in \mathcal{K}_-$ so that $[g,g]_{\mathcal{K}} < 0$. Unless the
observable probability functional is literally equal to the Krein inner
product, $[g,g]_{\mathcal{K}} < 0$ does not imply $P_{\mathrm{obs}} < 0$.
\end{proposition}

\begin{proof}
Observable probability is a functional on $M_{\mathrm{obs}}$, not on
$M_{\mathrm{adm}}$. It has the form $P_{\mathrm{obs}}(F \mid I) =
f(\Pi(g))$ for some function $f$ on $M_{\mathrm{obs}}$. The statement
$[g,g]_{\mathcal{K}} < 0$ is a statement about the Krein inner product on
$M_{\mathrm{adm}}$. Unless $f(\Pi(\cdot)) = [\cdot,\cdot]_{\mathcal{K}}$
identically --- which is precisely the Hilbert Assumption in Krein form
--- the sign of $[g,g]_{\mathcal{K}}$ is independent of the sign of
$P_{\mathrm{obs}}$.
\end{proof}

\subsection*{E.3 Ostrogradsky: State Pathology versus History Pathology}

\begin{proposition}
Ostrogradsky unboundedness is a pathology of the extended phase-space
representation, not of projected observable histories.
\end{proposition}

\begin{proof}
Ostrogradsky's theorem establishes
\begin{equation*}
  \inf_{(Q,P) \in T^*\!M_{\mathrm{adm}}^{\mathrm{ext}}} H(Q,P) = -\infty.
\end{equation*}
Observable stability is the assertion that projected histories remain
admissible:
\begin{equation*}
  A_O(\Pi(\gamma(t))) = 1 \quad \text{for all } t \in [0,T].
\end{equation*}
The first statement quantifies over all points in the extended phase space.
The second quantifies over actual dynamical trajectories and their
projections. A trajectory $\gamma(t)$ may avoid the unbounded regions of
$T^*\!M_{\mathrm{adm}}^{\mathrm{ext}}$ entirely --- the Pais--Uhlenbeck
oscillator does exactly this, as its solutions are bounded periodic orbits
despite the Hamiltonian being unbounded below. Therefore $H$ unbounded
below does not imply $\Pi(\gamma(t))$ unstable.
\end{proof}

\subsection*{E.4 Born Rule as a Projection from Rotation-Scaling Transport}

\begin{proposition}
The Born rule is a projection from the amplitude space (rotation-scaling
geometry) to the observable intensity space.
\end{proposition}

\begin{proof}
Write a quantum transition amplitude $a_{FI} = r_{FI}e^{i\theta_{FI}} \in \mathbb{C}$.
The Born rule assigns probability $P(F \mid I) = |a_{FI}|^2 = r_{FI}^2$.
The map $a_{FI} \mapsto |a_{FI}|^2$ is many-to-one: any two amplitudes
$re^{i\theta_1}$ and $re^{i\theta_2}$ with $\theta_1 \neq \theta_2$ produce
the same probability $r^2$. The phase $\theta_{FI}$ lies in the fiber of
this map; it is a hidden direction in amplitude space that does not appear
in $M_{\mathrm{obs}}$. The Born rule is therefore a surjective map from the
complex amplitude geometry of $M_{\mathrm{adm}}$ to the non-negative real
intensity geometry of $M_{\mathrm{obs}}$: a projection in the precise sense
of Definition~2.1.
\end{proof}

\subsection*{E.5 Irreversibility from Non-Invertibility of Projection}

\begin{theorem}[Irreversibility from Non-Invertibility]
If $\Pi : M_{\mathrm{adm}} \to M_{\mathrm{obs}}$ is non-injective, then
no reconstruction map $R : M_{\mathrm{obs}} \to M_{\mathrm{adm}}$ satisfies
$R \circ \Pi = \mathrm{id}_{M_{\mathrm{adm}}}$ on all of $M_{\mathrm{adm}}$.
\end{theorem}

\begin{proof}
Since $\Pi$ is non-injective, there exist $h \neq h'$ in $M_{\mathrm{adm}}$
with $\Pi(h) = \Pi(h')$. If $R \circ \Pi = \mathrm{id}$, then
$h = R(\Pi(h)) = R(\Pi(h')) = h'$, contradicting $h \neq h'$. Hence no
perfect reconstruction exists.
\end{proof}

\begin{definition}[Reconstruction Defect and Irreversibility]
Given a reconstruction map $R : M_{\mathrm{obs}} \to M_{\mathrm{adm}}$
and a metric $d$ on $M_{\mathrm{adm}}$, the \emph{reconstruction defect}
of $h$ is
\begin{equation}
  \Delta_R(h) = d\!\bigl(h,\, R(\Pi(h))\bigr).
\end{equation}
A system is \emph{irreversible at $h$} if $\Delta_R(h) > 0$ for every
reconstruction map $R$.
\end{definition}

\begin{remark}
Under this definition, irreversibility is not a property of a dynamics but
of a projection: it is the persistent failure of any map from observations
back to the generating history. Thermodynamic irreversibility, memory loss,
and the arrow of time all become special cases of the same structure ---
non-invertible projection --- distinguished only by the space $M_{\mathrm{adm}}$
and metric $d$ in question. This connects the admissibility framework to a
trajectory-based account of the second law.
\end{remark}

\subsection*{E.6 Jensen--Shannon Distance: Interface-Nativity and Fisher Expansion}

The following consolidates the properties of $d_{\mathrm{int}}$ established
in Section~8.2 into a single self-contained block.

\begin{proposition}[Interface-Nativity]
The Jensen--Shannon divergence $\operatorname{JSD}(P \| Q)$ compares
observable interfaces without reference to any underlying admissibility
manifold.
\end{proposition}

\begin{proof}
The definition $\operatorname{JSD}(P \| Q) = H(\frac{P+Q}{2}) - \frac{1}{2}(H(P) +
H(Q))$ uses only $P$, $Q$, their mixture $M = \frac{1}{2}(P+Q)$, and the
Shannon entropy $H$. No element of $M_{\mathrm{adm}}$ appears. The
computation is therefore intrinsic to the observable distribution space.
\end{proof}

\begin{theorem}[Gauge Equivalence as Zero Interface Distance]
Two histories $h, h' \in M_{\mathrm{adm}}$ are observationally equivalent
if and only if $\operatorname{JSD}(\Pi(h), \Pi(h')) = 0$.
\end{theorem}

\begin{proof}
Since $\operatorname{JSD}$ is a metric on probability distributions,
$\operatorname{JSD}(P, Q) = 0 \iff P = Q$. Therefore
$\operatorname{JSD}(\Pi(h), \Pi(h')) = 0 \iff \Pi(h) = \Pi(h'))$,
which is precisely the condition that $h$ and $h'$ lie in the same fiber
of $\Pi$ (Definition~8.1). Gauge orbits are therefore the zero-distance
fibers of the interface metric.
\end{proof}

\begin{proposition}[Fisher Expansion of Interface Metric]
For a parametric family $P_\lambda$ of observable distributions,
\begin{equation}
  \operatorname{JSD}^2(P_\lambda, P_{\lambda + \epsilon})
  = \frac{1}{8} \sum_x \frac{(\delta P(x))^2}{P_\lambda(x)} + O(|\epsilon|^3)
  = \frac{1}{8}\, g^F_{ij}(\lambda)\,\epsilon^i \epsilon^j + O(|\epsilon|^3),
\end{equation}
where $g^F_{ij}(\lambda) = \sum_x P_\lambda(x)\,\partial_i \log P_\lambda(x)\,
\partial_j \log P_\lambda(x)$ is the Fisher information metric.
\end{proposition}

\begin{proof}
Let $\delta P(x) = P_{\lambda+\epsilon}(x) - P_\lambda(x)$ and
$M(x) = P_\lambda(x) + \frac{1}{2}\delta P(x)$. For small $\epsilon$,
the KL expansion $D_{\mathrm{KL}}(P \| M) = \sum_x P(x)\log(P(x)/M(x))$
yields $D_{\mathrm{KL}}(P \| M) = \frac{1}{8}\sum_x (\delta P(x))^2/P_\lambda(x)
+ O(|\epsilon|^3)$. Since $\operatorname{JSD}(P\|Q) = \frac{1}{2}D_{\mathrm{KL}}(P\|M)
+ \frac{1}{2}D_{\mathrm{KL}}(Q\|M)$, and both KL terms contribute equally
at leading order, the result follows with $\delta P(x) = \partial_i P_\lambda(x)\,\epsilon^i$.
\end{proof}

\subsection*{E.7 Distinction Loss Functional}

\begin{definition}[Distinction Functional]
For observable distributions $P, Q$ on $M_{\mathrm{obs}}$, the
\emph{distinction functional} is
\begin{equation}
  \mathcal{D}(P, Q)
  = H\!\left(\frac{P+Q}{2}\right) - \frac{1}{2}\bigl(H(P) + H(Q)\bigr)
  = \operatorname{JSD}(P \| Q).
\end{equation}
\end{definition}

\begin{proposition}
$\mathcal{D}(P,Q)$ measures the information destroyed when the distinction
between $P$ and $Q$ is erased by replacing both with their mixture.
\end{proposition}

\begin{proof}
The entropy $H(\frac{P+Q}{2})$ is the uncertainty of an observer who
cannot distinguish which interface generated each observation. The average
$\frac{1}{2}(H(P) + H(Q))$ is the uncertainty when the distinction is
retained. Their difference $\mathcal{D}(P,Q) \geq 0$ is the additional
uncertainty caused by forgetting the distinction. It vanishes if and only
if $P = Q$, in which case the distinction carries no information.
\end{proof}

\subsection*{E.8 Summary Correspondence Table}

The following table translates every major concept of the admissibility
framework into the language of interface geometry and JSD.

\medskip
{\small
\begin{center}
\begin{tabular}{p{3.8cm}p{5.8cm}}
\toprule
Admissibility concept & Interface-geometric translation \\
\midrule
Gauge orbit &
Zero-JSD fiber of $\Pi$: $\operatorname{JSD}(\Pi(h), \Pi(h')) = 0$ \\
Hidden admissibility direction &
Deformation with vanishing first-order $\delta_{\mathrm{obs}}$ \\
Ghost state &
Element of $\mathcal{K}_-$ mapping to zero-JSD deformation \\
Hilbert Assumption &
$\Pi$ injective: no non-trivial zero-JSD fibers \\
Observable pathology &
$\mathcal{D}(P, P') > 0$ induced by the deformation \\
Irreversibility &
Positive reconstruction defect $\Delta_R > 0$ \\
Born rule &
Projection $re^{i\theta} \mapsto r^2$ from amplitude to intensity \\
Theory comparison &
$d_{\mathrm{int}}(P, Q) = \operatorname{JSD}^{1/2}(P \| Q)$ \\
Fisher geometry of theories &
$g_{ij} = \frac{1}{8} I_{ij}(\lambda)$ at leading order \\
Derived Interface Conjecture &
$d_{\mathrm{int}}(\text{QG},\, \text{quadratic gravity}) = 0$ below $M_{\mathrm{Pl}}$ \\
\bottomrule
\end{tabular}
\end{center}}
\medskip

\subsection*{E.9 The Mature Framework in One Chain}

The progression from internal machinery to distinguishability geometry
can be compressed into a single diagram:
\begin{equation}
  M_{\mathrm{adm}}
  \;\xrightarrow{\;\Pi\;}
  M_{\mathrm{obs}}
  \;\xrightarrow{\;\operatorname{JSD}\;}
  (\mathcal{M},\, d_{\mathrm{int}}).
\end{equation}
The first arrow discards internal structure; what survives is the observable
interface. The second arrow measures distinguishability between interfaces;
what survives is the geometry of distinction. The question
\begin{equation}
  \text{``What exists?''}
\end{equation}
is replaced by
\begin{equation}
  \text{``What distinctions survive projection?''}
\end{equation}
Ghost states, gauge orbits, and Hilbert-space structure are all answers to
the question of what does \emph{not} survive: they are zero-JSD fibers,
hidden directions, and special cases of injective projection, respectively.
The geometry of what does survive --- the metric space
$(\mathcal{M}, d_{\mathrm{int}})$ --- is the object that the theory of
physical representation is ultimately about.

\begin{thebibliography}{99}

\bibitem{stelle1977}
K.~S. Stelle,
``Renormalization of higher-derivative quantum gravity,''
\textit{Physical Review D} \textbf{16}, 953--969 (1977).

\bibitem{stelle1978}
K.~S. Stelle,
``Classical gravity with higher derivatives,''
\textit{General Relativity and Gravitation} \textbf{9}(4), 353--371 (1978).

\bibitem{avramidi1985}
I.~G. Avramidi and A.~O. Barvinsky,
``Asymptotic freedom in higher-derivative quantum gravity,''
\textit{Physics Letters B} \textbf{159}(4--6), 269--274 (1985).

\bibitem{pais1950}
A.~Pais and G.~E. Uhlenbeck,
``On field theories with non-localized action,''
\textit{Physical Review} \textbf{79}, 145--165 (1950).

\bibitem{bender2008}
C.~M. Bender and P.~D. Mannheim,
``No-ghost theorem for the fourth-order derivative Pais--Uhlenbeck oscillator model,''
\textit{Physical Review Letters} \textbf{100}, 110402 (2008).

\bibitem{ostrogradsky1850}
M.~V. Ostrogradsky,
``M\'{e}moires sur les \'{e}quations diff\'{e}rentielles,
relatives au probl\`{e}me des isop\'{e}rim\`{e}tres,''
\textit{M\'{e}moires de l'Acad\'{e}mie imp\'{e}riale des sciences
de Saint-P\'{e}tersbourg} \textbf{VI}(4), 385--517 (1850).

\bibitem{woodard2015}
R.~P. Woodard,
``Ostrogradsky's theorem on Hamiltonian instability,''
\textit{Scholarpedia} \textbf{10}(8), 32243 (2015).

\bibitem{mannheim2012}
P.~D. Mannheim,
``Making the case for conformal gravity,''
\textit{Foundations of Physics} \textbf{42}, 388--420 (2012).

\bibitem{salvio2018}
A.~Salvio and A.~Strumia,
``Agravity up to infinite energy,''
\textit{European Physical Journal C} \textbf{78}, 124 (2018).

\bibitem{anselmi2017}
D.~Anselmi and M.~Piva,
``A new formulation of Lee--Wick quantum field theory,''
\textit{Journal of High Energy Physics} \textbf{2017}(6), 066 (2017).

\bibitem{krein1978}
M.~G. Kre\u{\i}n and H.~Langer,
``On some mathematical principles in the linear theory of damped oscillations
of continua,''
\textit{Integral Equations and Operator Theory} \textbf{1}(3), 364--399 (1978).

\bibitem{bleuler1950}
K.~Bleuler,
``Eine neue Methode zur Behandlung der longitudinalen und skalaren Photonen,''
\textit{Helvetica Physica Acta} \textbf{23}, 567--586 (1950).

\bibitem{faddeev1967}
L.~D. Faddeev and V.~N. Popov,
``Feynman diagrams for the Yang--Mills field,''
\textit{Physics Letters B} \textbf{25}(1), 29--30 (1967).

\bibitem{becchi1976}
C.~Becchi, A.~Rouet, and R.~Stora,
``Renormalization of gauge theories,''
\textit{Annals of Physics} \textbf{98}(2), 287--321 (1976).

\bibitem{turok2024}
N.~Turok and L.~Boyle,
``A minimal explanation of the primordial cosmological perturbations,''
\textit{Physical Review Letters} \textbf{132}, 061402 (2024).

\bibitem{birrell1982}
N.~D. Birrell and P.~C.~W. Davies,
\textit{Quantum Fields in Curved Space},
Cambridge University Press, Cambridge, 1982.

\bibitem{donoghue1994}
J.~F. Donoghue,
``General relativity as an effective field theory: the leading quantum
corrections,''
\textit{Physical Review D} \textbf{50}, 3874--3888 (1994).

\bibitem{weinberg1979}
S.~Weinberg,
``Ultraviolet divergences in quantum theories of gravitation,''
in \textit{General Relativity: An Einstein Centenary Survey},
ed.\ S.~W.~Hawking and W.~Israel,
Cambridge University Press, Cambridge, 1979, pp.~790--831.

\bibitem{deser2009}
S.~Deser,
``Ostrogradski, ghosts, and gauge as shadows of higher derivatives,''
talk given at the International School on Strings and Fundamental Physics,
Munich, 2009.

\bibitem{needham1997}
T.~Needham,
\textit{Visual Complex Analysis},
Oxford University Press, Oxford, 1997.

\bibitem{barandes2023}
J.~A. Barandes,
``The unistochastic quantum theory,''
\textit{arXiv:2302.10778} [quant-ph] (2023).

\bibitem{barandes2024}
J.~A. Barandes,
``Is quantum mechanics just thermodynamics in disguise?''
\textit{arXiv:2407.09663} [quant-ph] (2024).

\bibitem{endres2003}
D.~M. Endres and J.~E. Schindelin,
``A new metric for probability distributions,''
\textit{IEEE Transactions on Information Theory} \textbf{49}(7),
1858--1860 (2003).

\bibitem{amari2000}
S.-I. Amari and H. Nagaoka,
\textit{Methods of Information Geometry},
American Mathematical Society, Providence, 2000.



\bibitem{boyle2022dim}
L.~Boyle and N.~Turok,
``Cancelling the vacuum energy and Weyl anomaly in the standard model
with dimension-zero scalar fields,''
\textit{arXiv:2110.06258} [hep-th] (2022).

\bibitem{boyle2018cpt}
L.~Boyle, K.~Finn, and N.~Turok,
``CPT-symmetric universe,''
\textit{Physical Review Letters} \textbf{121}, 251301 (2018).

\bibitem{boylefinn2022}
L.~Boyle, K.~Finn, and N.~Turok,
``The Big Bang, CPT, and neutrino dark matter,''
\textit{Annals of Physics} \textbf{438}, 168767 (2022).

\bibitem{turokboyle2022entropy}
N.~Turok and L.~Boyle,
``Gravitational entropy and the flatness, homogeneity and isotropy puzzles,''
\textit{Physics Letters B} \textbf{849}, 138442 (2024).

\bibitem{clinebell2025}
J.~M. Cline and A.~Hell,
``Pathologies of dimension-zero scalar fields,''
\textit{arXiv:2603.05683} [hep-th] (2025).

\bibitem{boyle2023perturbations}
L.~Boyle and N.~Turok,
``Thermodynamic solution of the homogeneity, isotropy and flatness puzzles
(and a clue to the cosmological constant),''
\textit{Physics Letters B} \textbf{849}, 138443 (2024).

\bibitem{bender2008b}
C.~M. Bender and P.~D. Mannheim,
``Exactly solvable $\mathcal{PT}$-symmetric models in two dimensions,''
\textit{Journal of Physics A: Mathematical and Theoretical} \textbf{41}, 244005 (2008).

\bibitem{bogolubov1990}
N.~N. Bogolubov, A.~A. Logunov, A.~I. Oksak, and I.~T. Todorov,
\textit{General Principles of Quantum Field Theory},
Kluwer Academic Publishers, Dordrecht, 1990.

\bibitem{hawkinghertog2002}
S.~W. Hawking and T.~Hertog,
``Living with ghosts,''
\textit{Physical Review D} \textbf{65}, 103515 (2002).

\bibitem{duff1994}
M.~J. Duff,
``Twenty years of the Weyl anomaly,''
\textit{Classical and Quantum Gravity} \textbf{11}, 1387--1404 (1994).

\bibitem{penrose1964}
R.~Penrose,
``Conformal treatment of infinity,''
in \textit{Relativity, Groups and Topology},
eds.\ C.~DeWitt and B.~S.~DeWitt,
Gordon and Breach, London, 1964.

\end{thebibliography}

\end{document}

